%!TEX TS-program = xelatex
%!TEX encoding = UTF-8 Unicode
\documentclass[manuscript,screen,review]{acmart}

\setcopyright{none}
\copyrightyear{2026}
\acmYear{2026}
\acmJournal{TOSEM}
\settopmatter{printacmref=true, printccs=true, printfolios=true}
\setcitestyle{numbers,sort&compress,square}

\usepackage{amssymb,amsthm}
\usepackage{mathtools}
\usepackage{calc}
\usepackage{enumitem}
\usepackage{longtable}
\usepackage{array}
\usepackage{listings}
\usepackage{pifont}

\setlength{\emergencystretch}{3em}
\lstset{
  breaklines=true,
  breakatwhitespace=true,
  basicstyle=\ttfamily\small,
  frame=none,
  columns=fullflexible
}

\title[Supplementary Material]{Supplementary Material for ``A Semantic Mutation Metric for Metamorphic-Relation Adequacy in Scientific Computing Programs''}

\author{Meng Li}
\orcid{0000-0002-1074-1502}
\email{mlemon@usc.edu.cn}
\affiliation{%
  \institution{School of Computing, University of South China}
  \city{Hengyang}
  \country{China}
}
\affiliation{%
  \institution{Hunan Engineering Research Center of Software Evaluation and Testing for Intellectual Equipment}
  \city{Hengyang}
  \country{China}
}
\affiliation{%
  \institution{CNNC Key Laboratory on High Trusted Computing}
  \city{Hengyang}
  \country{China}
}

\author{Xiaohua Yang}
\orcid{0000-0002-2977-1787}
\email{xiaohua1963@foxmail.com}
\affiliation{%
  \institution{School of Computing, University of South China}
  \city{Hengyang}
  \country{China}
}

\author{Jie Liu}
\orcid{0009-0008-1970-8347}
\email{jieliu@usc.edu.cn}
\affiliation{%
  \institution{School of Computing, University of South China}
  \city{Hengyang}
  \country{China}
}

\author{Shiyu Yan}
\orcid{0000-0001-7626-5185}
\email{yanshiyu@usc.edu.cn}
\affiliation{%
  \institution{School of Computing, University of South China}
  \city{Hengyang}
  \country{China}
}

\begin{CCSXML}
<ccs2012>
 <concept>
  <concept_id>10011007.10011074.10011099.10011102</concept_id>
  <concept_desc>Software and its engineering~Software testing and debugging</concept_desc>
  <concept_significance>500</concept_significance>
 </concept>
 <concept>
  <concept_id>10011007.10011074.10011099</concept_id>
  <concept_desc>Software and its engineering~Software verification and validation</concept_desc>
  <concept_significance>300</concept_significance>
 </concept>
</ccs2012>
\end{CCSXML}

\ccsdesc[500]{Software and its engineering~Software testing and debugging}
\ccsdesc[300]{Software and its engineering~Software verification and validation}

\keywords{metamorphic testing, mutation testing, semantic mutation testing, metamorphic relation adequacy, metamorphic relation assessment, scientific computing}

\begin{document}

\maketitle

\appendix

\section{Notation and Operator Catalogue}\label{a.-notation-and-operator-catalogue}

\subsection{Complete notation table}\label{a.1-complete-notation-table}

\paragraph{Experimental objects.}
\begin{itemize}
\item PUT set $I = \{\text{A1,A2,A3, B1,B2,B3, C1,C2,C3, D1,D2,D3}\}$, $|I| = 12$.
\item PUT: $S_i$ ($i \in I$).
\item Class mapping: $\mathrm{cls} : I \to \{A, B, C, D\}$ (open, extensible).
\item Valid input distribution: $\mathcal{D}_P$, sampling $X_{K_{\mathrm{eq}}} \sim \mathcal{D}_P$.
\end{itemize}

\paragraph{Metamorphic testing side.}
\begin{itemize}
\item Metamorphic-relation label set: $\mathcal{R} = \{\text{Cons.}, \text{Mono.}, \text{Conv.}, \text{Traj.}, \text{P-ord.}\}$, $k \in \{1,\ldots,5\}$ (open, extensible).
\item Metamorphic Relation: $\mathrm{MR}_{i,k}$ (from the shared infrastructure); $\mathit{mr} = (r, R) \in \mathrm{MR}_{i,k}$.
\item Automated Verification Pipeline (AVP): $\mathrm{AVP} : \text{Programs} \times \mathrm{MR}_{\text{universe}} \times \mathbb{R}^+ \to \{\text{pass}, \text{fail}\}$; $\mathrm{AVP}(s, \mathit{mr}, \varepsilon_{\mathrm{AVP}}^k)$ invokes the verification method per relation index $k$. Where the tolerance is fixed by the relation's meta-pattern we abbreviate $\mathrm{AVP}(s, \mathit{mr}) := \mathrm{AVP}(s, \mathit{mr}, \varepsilon_{\mathrm{AVP}}^k)$.
\end{itemize}

\paragraph{Mutation testing side.}
\begin{itemize}
\item Mutation operator family: $\mathrm{MUT} = \{\mathrm{mut}_1, \ldots, \mathrm{mut}_5\}$, $j \in \{1,\ldots,5\}$ (open, extensible); $\mathrm{mut}_1 = \mathrm{mut}_C$, $\mathrm{mut}_2 = \mathrm{mut}_M$, $\mathrm{mut}_3 = \mathrm{mut}_G$, $\mathrm{mut}_4 = \mathrm{mut}_T$, $\mathrm{mut}_5 = \mathrm{mut}_F$.
\item Signature: $\mathrm{mut}_j : \text{Programs} \to 2^{\text{Programs}}$.
\item Mutant: $s' \in \mathrm{mut}_j(S_i)$.
\item Alignment: $\mathrm{align}(j) = j$.
\end{itemize}

\paragraph{Evaluation unit and pool decomposition (fixed).}
The evaluation unit is the (PUT, meta-pattern) cell $(i,k)$ over the
PUT's pooled admitted mutants
$\mathrm{mut}(S_i) = \bigcup_j \mathrm{mut}_j(S_i)$; the operator index
$j$ is a per-mutant fiber label, not a cell axis.
\[
\begin{aligned}
\mathrm{mut}(S_i)
  &= \mathrm{certified}_{i,k} \cup \mathrm{killed}_{i,k}
     \cup \mathrm{survive}_{i,k} \cup \mathrm{unresolved}_{i,k},\\
&\hspace{1.5em}\text{with the four sets disjoint, mutually exclusive, and exhaustive.}
\end{aligned}
\]
\begin{itemize}
\item \emph{equivalence-status evidence} (dual conditions):
  (E1) AVP-coherent: $\forall \mathit{mr} \in \mathrm{MR}_{i,k}: \mathrm{AVP}(S_i, \mathit{mr}) = \mathrm{AVP}(s', \mathit{mr})$;
  (E2) Output-equiv: $\forall x \in X_{K_{\mathrm{eq}}} \sim \mathcal{D}_P: \|S_i(x) - s'(x)\| \leq \varepsilon_{\mathrm{eq}}$.
  E1 $\vee$ E2 failure yields \texttt{CONFIRMED\_NON\_EQUIVALENT}; E1
  $\wedge$ E2 agreement without a machine-checkable certificate yields
  \texttt{EQUIVALENCE\_UNRESOLVED}; \texttt{CERTIFIED\_EQUIVALENT} is
  reserved for certificate-bearing candidates (empty in the current
  pipeline).
\item \emph{killed determination}: $\mathrm{killed}(s', \mathrm{MR}_{i,k}) \Leftrightarrow \exists \mathit{mr} \in \mathrm{MR}_{i,k}: \mathrm{AVP}(S_i, \mathit{mr}) = \text{pass} \wedge \mathrm{AVP}(s', \mathit{mr}) = \text{fail}$; killed mutants are confirmed non-equivalent (kill witness upgrade, main-text lemma).
\end{itemize}

\paragraph{Metric (fixed classical structure; strict and conservative readings).}
\begin{align*}
\mathrm{SMS}_{i,k} := \mathrm{SMS}^{\mathrm{strict}}_{i,k} &= \frac{|\mathrm{killed}_{i,k}|}{|\mathrm{killed}_{i,k}| + |\mathrm{survive}_{i,k}|}, \\
\mathrm{SMS}^{\mathrm{cons}}_{i,k} &= \frac{|\mathrm{killed}_{i,k}|}{|\mathrm{killed}_{i,k}| + |\mathrm{survive}_{i,k}| + |\mathrm{unresolved}_{i,k}|} \in [0, 1],
\end{align*}
each \texttt{NA} on an empty denominator. With $u_{i,k} = 0$ (the
audited v4 campaign) both equal the classical ratio
$|\mathrm{killed}|/(|\mathrm{mut}| - |\mathrm{equiv}|)$.

\paragraph{LRCA engineering attribution layer (descriptive, not in SMS).}
\begin{itemize}
\item Likely root-cause inventory: $C = \{C_1, \ldots, C_5\}$ (open, extensible). C1 Likely semantic failure (residual label: no C2--C5 trigger); C2 Tolerance perturbation; C3 OOD; C4 Statistical-assumption violation; C5 Mutator artefact.
\item LRCA output: each $s' \in \mathrm{killed}$ annotated with $\texttt{root\_cause}(s') \in C$.
\item Descriptive quantities: $\texttt{C1\_share}_{i,k}$, $\texttt{suspect\_share}_{i,k}$; both \texttt{NA} on zero-kill cells.
\end{itemize}

\paragraph{Slicing and cross-class.}
\begin{itemize}
\item Slicing: aligned $k = k^\star(i)$ (the PUT's class-primary meta-pattern); cross $k \neq k^\star(i)$.
\item Cross-class: $\Delta\mathrm{SMS}_c$, $c \in \{A, B, C, D\}$; $\mathrm{CV}(\Delta\mathrm{SMS}) := \mathrm{std}(\Delta\mathrm{SMS}_c) / |\mathrm{mean}(\Delta\mathrm{SMS}_c)|$.
\end{itemize}

\paragraph{Theory-extension symbols (window, interval, and gap results).}
\begin{itemize}
\item Observation: $\mathrm{obs} : \mathcal{Y} \to \mathcal{Y}_{\mathrm{obs}}$; observational equivalence $\equiv_{\mathrm{obs}}$; computation map $\Phi_Q$.
\item Invariant family: $\Psi = \{\psi_1, \dots, \psi_5\}$ (non-redundant; open, e.g.\ $\psi_6$); edit operator $\mathrm{edit}$ with $P_{\mathrm{edit}} = \mathrm{edit}(P)$; effect map $\mathrm{eff}$ with fibers $\mathrm{eff}^{-1}(\cdot)$.
\item Counts (local to the theory subsections): $n$ confirmed non-equivalent, $\kappa$ killed, $u$ unresolved, $u_{\mathrm{neq}} \in [0, u]$.
\item Coverage and gaps: $\mathrm{Cov}(R)$; $\mathrm{Gap}_{\mathrm{aln}}(R)$; $\mathrm{Gap}_{\mathrm{str}}(R)$; exactness defect $\xi(R)$ (\texttt{NA} on zero kills; one-sided); kill signature $\mathrm{sig}(m)$; universe $M_{\mathrm{neq}}$, fibers $M_j$, weights $w_j$.
\item Window quantities: violation magnitude $\varepsilon_{\mathrm{viol}}(m)$; residual $\Delta_r = \sup_{x \in D_r} \varepsilon_r(x; P^\star)$ over the applicable domain $D_r$; margin $\mu_r = \varepsilon_{\mathrm{tol}} - \Delta_r$; per-run noise bound $\bar\eta$ with execution multiplicity $p$; repeats $N$, concentration constant $c$, output spread $\sigma_{\mathrm{out}}$, deterministic noise $\eta_{\mathrm{det}}$; Lipschitz constant $L_r$; crash threshold $\varepsilon_{\mathrm{crash}}$; exception set $\mathcal{N}_{\mathrm{exc}}$ (degeneration).
\item Degenerate limit: $L = L_{\mathrm{lim}} \wedge L_{\mathrm{switch}} = L_{\mathrm{equiv}} \wedge L_{\mathrm{killed}} \wedge L_{\mathrm{mut}}$.
\end{itemize}

\paragraph{Tolerance and sampling.}
\begin{itemize}
\item Tolerance: $\varepsilon_{\mathrm{AVP}}^k$ (relation-dependent), $\varepsilon_{\mathrm{eq}}$ (equivalence output tolerance).
\item Sampling: $K_{\mathrm{eq}} = 1000$ (E2 sample size), $N = 20$ (statistical replicates).
\end{itemize}

The notation skeleton (11 core concepts + three-state decomposition +
SMS formula + AVP interface) is \textbf{fixed}; the \emph{content} of
MUT, $\mathcal{R}$, C, and cls is open to extension.

\subsection{Root-Cause Decision Tree and Diagnostic Protocol}\label{a.2-lrca-decision-tree-and-three-layer-diagnostic-protocol}

\textbf{Five likely root causes} (open):

{%
\begin{longtable}[]{@{}ll@{}}
\caption{Root-cause decision tree and three-layer diagnostic protocol.}\label{tab:p2-15}\\
\toprule\noalign{}
Code & Meaning \\
\midrule\noalign{}
\endfirsthead
\endhead
\bottomrule\noalign{}
\endlastfoot
C1 & Likely semantic failure (residual label: no C2--C5 trigger; the
kill class SMS seeks) \\
C2 & Numerical tolerance perturbation \\
C3 & Out-of-distribution (OOD) \\
C4 & Statistical-assumption violation \\
C5 & Mutator artefact \\
\end{longtable}
}

\textbf{Three-layer diagnosis with L0 artefact pre-scan:}

\begin{itemize}
\item \textbf{L0 artefact pre-scan.} Double-blind review following the main manuscript's mutant-pool prescreen and equivalence-detection protocol: dual-LLM cross-source review plus 20\% manual sampling.
\item \textbf{L1 tolerance robustness.} $N = 20$ replicates; fail ratio $\geq 0.80$ considered stable.
\item \textbf{L2 OOD triage.} For C/D classes, distinguish valid $\mathcal{D}_P^{\text{valid}}$ from OOD.
\item \textbf{L3 statistical baseline.} For Wilcoxon/DTW, IID/stationarity pre-check on PUT samples.
\end{itemize}

\textbf{Decision tree:}

For each $s' \in \mathrm{killed}_{i,k}$:
\begin{itemize}
\item \textbf{L1.} fail ratio $< 0.80 \Rightarrow$ C2; otherwise $\Rightarrow$ L2.
\item \textbf{L2} (C/D classes). fail only in OOD $\Rightarrow$ C3; otherwise $\Rightarrow$ L3.
\item \textbf{L3} (B/D classes + Wilcoxon/DTW). assumption violation $\Rightarrow$ C4; otherwise $\Rightarrow$ artefact recheck.
\item \textbf{Artefact recheck.} artefact evidence $\Rightarrow$ C5; otherwise $\Rightarrow$ C1.
\end{itemize}

\textbf{Multi-label priority:} C5 \textgreater{} C4 \textgreater{} C3
\textgreater{} C2 \textgreater{} C1.

\textbf{LRCA threshold calibration (9-grid scan;
\texttt{lrca\_calibration.json}).} Best ood\_band=0.02 with any
tolerance\_multiplier (3.0 / 10.0 / 30.0) gives 12 of 15 evaluable
cells at suspect\_share $\leq 0.20$.
\texttt{tolerance\_multiplier} has zero impact (L1 tolerance
rarely triggered); \texttt{ood\_band} is the only discriminator. Under
the \texttt{NA} convention of the main text the pass-count denominator
is the 15 evaluable cells, not 60; the main-text H4 discussion renders
no pass/fail verdict because the pre-registered 60-cell mean is
undefined on this data. The grid calibrates diagnostic labels only and
is not an H4 hypothesis test.

\subsection{Verification Protocol and Equivalence-Judgement Trade-Off}\label{a.3-avp-protocol-specification-and-equivalence-judgement-trade-off}

\[
\mathrm{AVP}: \mathrm{Programs} \times \mathrm{MR}_{\mathrm{universe}} \times \mathbb{R}^{+} \to \{\mathrm{pass}, \mathrm{fail}\}.
\]
Per-relation verification: $\mathrm{Cons.}$ tolerance equality $|\mathrm{LHS} -
\mathrm{RHS}| \leq \varepsilon$; $\mathrm{Mono.}$ / $\mathrm{P-ord.}$ Wilcoxon signed-rank $\alpha=0.05$; $\mathrm{Conv.}$
convergence-order + asymptotic residual ratio; Traj. DTW distance
threshold $\varepsilon_{\mathrm{DTW}}$. AVP version = upstream infrastructure commit hash; our package embeds
AVP source.

\textbf{Three-candidate equivalence trade-off.}

{%
\begin{longtable}[]{@{}
  >{\raggedright\arraybackslash}p{(\linewidth - 6\tabcolsep) * \real{0.2500}}
  >{\raggedright\arraybackslash}p{(\linewidth - 6\tabcolsep) * \real{0.2500}}
  >{\raggedright\arraybackslash}p{(\linewidth - 6\tabcolsep) * \real{0.2500}}
  >{\raggedright\arraybackslash}p{(\linewidth - 6\tabcolsep) * \real{0.2500}}@{}}
\caption{Verification protocol and equivalence-judgement trade-off.}\label{tab:p2-16}\\
\toprule\noalign{}
\begin{minipage}[b]{\linewidth}\raggedright
Determination
\end{minipage} & \begin{minipage}[b]{\linewidth}\raggedright
False-positive
\end{minipage} & \begin{minipage}[b]{\linewidth}\raggedright
False-negative
\end{minipage} & \begin{minipage}[b]{\linewidth}\raggedright
SMS bias
\end{minipage} \\
\midrule\noalign{}
\endfirsthead
\endhead
\bottomrule\noalign{}
\endlastfoot
E1 alone & numerical coincidence + AVP-coherent & E1 false $\to$ non-equiv &
Low (smaller denominator) \\
E2 alone & output match, MR differs (rare) & K\_eq miss, full-space
equal & High \\
\textbf{E1 $\wedge$ E2} & both mis-pass simultaneously & E1 $\vee$ E2 false $\to$ non-equiv
& Slightly low (conservative: misclassified equivalents remain as
unkillable survivors in the denominator) \\
\end{longtable}
}

E1 alone is fooled by insufficient AVP coverage; E2 alone by numerical
coincidence on $K_{\mathrm{eq}}$. E1 $\wedge$ E2 is the conservative complete
instantiation; under $L_{\mathrm{equiv}}$ (with the switch component
fixed) it reduces to
classical bitwise equivalence outside the exception set
$\mathcal{N}_{\mathrm{exc}}$ (Lemma G.1). Counter-examples: E2 passes
but E1 does not (rare; mutant coincides at K\_eq but deviates at MR
trigger points $\to$ non-equivalent); E1 passes but E2 does not (common;
consistent within MR but numerical drift \textgreater{} $\varepsilon_{\mathrm{eq}} \to$
non-equivalent).

% [thematic break removed]

\section{Experimental Subjects and Mutation-Operator Specialisations}\label{b.-experimental-subjects-and-mutation-operator-specialisations}

\subsection{Program Selection Rationale}\label{b.1-put-selection-rationale-detailed-coverage-argument}

\textbf{(a) Library-stack coverage.} numpy 2.4.4 (A1-B3 linear algebra /
arrays), scipy 1.17.1 (A1 integrate, A2 linalg, B2 stats), scikit-learn
1.8.0 (C1-D3 surrogate / ML), Python scientific-computing's de facto
foundation stack (PyPI top-50, 2026-04). Uncovered: GPU / distributed
(JAX, CuPy, dask), domain-specific (BioPython, Astropy, RDKit). Left to
future work under the representativeness threat.

\textbf{(b) Mathematical-structure coverage.} ODE/PDE (A1, A3), direct
linear algebra (A2), Bayesian analytic + MCMC + Monte Carlo (B1-B3),
kernel + orthogonal polynomial + NN surrogate (C1-C3), convex
optimisation + backprop + max-likelihood (D1-D3), covers 8 of 12
chapters in Numerical Recipes \citep{press2007numerical}. Uncovered: advanced
PDE solvers (FEM, FV, spectral); FFT; interior-point/trust-region
optimisation; symbolic/CAS. Left to future work.

\textbf{(c) Comparison with benchmarks.} DeepCrime (Humbatova et
al.~2021): semantic class-D topical overlap on ML kernels. Defects4J (Just et
al.~2014): no overlap, only mutation-as-fault-proxy reference. mutmut /
cosmic-ray demos: general Python; the semantic operators extend the scientific-computing focus.

\textbf{(d) Scale vs representativeness.} Each PUT kernel file is 19-35
LOC (0.6-1.2 KB), wrapping its library calls behind the signature
standardised to \texttt{program(x:\ float)\ $\to$\ float}. Substantively
smaller than industrial code (typically 1-10 KLOC).

\textbf{Limitation.} The \texttt{float\ $\to$\ float} signature is a
substantive constraint, not engineering trade-off. Industrial inputs are
typically high-dimensional (CFD grids, MD particles, FEM matrices);
scalarised PUTs upper-bound mutant semantic complexity and may
systematically under-estimate SMS and cross-class differences on
industrial PUTs (declared in the main manuscript's threats-to-validity section;
future work validates on industrial-grade PUTs).

\subsubsection{Systematic vs incidental}\label{b.1.5-systematic-vs-incidental}

Syntactic tools may \emph{occasionally} hit the semantic-mutation criteria in the main manuscript (a)
(cross-function-boundary) or (c) (algorithm class), but occasionality
undermines two engineering functions of semantic mutation. \textbf{(a)
Deepening source-code understanding.} Designing OS
\texttt{np.linalg.det(M)\ $\to$\ np.sum(np.diag(M))} requires knowing the
two are equivalent on diagonal matrices but not on general matrices
(\texttt{det\ =\ $\prod$\ eigvals} vs
\texttt{sum(diag)\ =\ trace\ =\ $\sum$\ eigvals}); syntactic tool AST
traversal does not require this understanding even when similar
replacements appear. \textbf{(b) Revealing deep faults.}
Domain-semantic errors (physical-constant, unit conversion, boundary
conditions, hyperparameter semantics, numerical-method order) are not
AST-local; syntactic mutator design goals (operator typos, off-by-one,
negation flips) have hit probability for domain errors much smaller than
systematic semantic-mutator triggering, and lack repeatability.
\textbf{Conclusion.} Syntactic tools occasionally producing mutants
satisfying (a)(b)(c) is stochastic byproduct, not repeatable, not
carrying engineering value. Systematic semantic mutation requires
(a)(b)(c) to be design intent.

\subsection{Per-class operator specialisations}\label{b.2-per-class-operator-specialisations}

{%
\begin{longtable}[]{@{}
  >{\raggedright\arraybackslash}p{(\linewidth - 2\tabcolsep) * \real{0.5000}}
  >{\raggedright\arraybackslash}p{(\linewidth - 2\tabcolsep) * \real{0.5000}}@{}}
\caption{Per-class operator specialisations.}\label{tab:p2-17}\\
\toprule\noalign{}
\begin{minipage}[b]{\linewidth}\raggedright
Operator
\end{minipage} & \begin{minipage}[b]{\linewidth}\raggedright
Representative specialisations across PUTs
\end{minipage} \\
\midrule\noalign{}
\endfirsthead
\endhead
\bottomrule\noalign{}
\endlastfoot
\textbf{mut\_C} Conservation-breaking & A1 Lorenz: add $\varepsilon_{\mathrm{drift}}$ to RHS
(slow Hamiltonian drift); A2 LU: decomposition omits k+1-th multiplier
(breaks det conservation); B1 Beta-Bin: posterior omits normalisation;
C1 GPR: covariance omits positive-definite diagonal term; D1 MLP:
backprop omits one gradient term. \\
\textbf{mut\_M} Monotonicity-breaking & A3 FDM: $\Delta$t coefficient
occasionally negative; B2 MCMC: acceptance min(1, r) $\to$ min(0.95, r); C2
PCE: high-order coefficient sort inserts inversion; D2 SVM:
decision-function sign flips near boundary. \\
\textbf{mut\_G} Convergence-breaking & A1 Lorenz: RK4 $\to$ 1.5-order
hybrid; A3 FDM: 2nd-order difference $\to$ 1st-order; B3 MC: doubling sample
size does not 1/N-reduce variance; C3 NN-Surr: training-epoch
truncation. \\
\textbf{mut\_T} Trajectory-distorting & A1 Lorenz: state-vector y/z
swap; B2 MCMC: insert independent-sampling segment; C3 NN-Surr:
training-target slow phase shift; D1 MLP: hidden-layer periodic-mask
activation. \\
\textbf{mut\_F} Fidelity-order-breaking & A2 LU: partial pivoting
degrades to no pivoting; C1 GPR: length-scale switches to coarse prior;
C2 PCE: high-order term randomly retains low-order; D3 LR:
regularisation occasionally large. \\
\end{longtable}
}

\textbf{Per-class HP / OS / TF substitution rules} (illustrative).
\textbf{HP} on class a: tolerance / max\_iter; on class c: GPR
\texttt{noise\_level} / \texttt{length\_scale}; on class d: MLP
\texttt{hidden\_dim} / dropout. \textbf{OS} on class a: numerical-linalg
API swaps (\texttt{det} $\leftrightarrow$ \texttt{sum(diag)}); on class b:
probability-distribution sampling swaps; on class c: surrogate-class
swaps (GPR $\leftrightarrow$ RBF $\leftrightarrow$ NN). \textbf{TF} on class a: integration-order
changes (RK4 $\to$ Euler); on class b: MC estimator changes.

\textbf{Operator-level cross-table with cosmic-ray defaults}. The 13
cosmic-ray default operators (the B.6 inventory) are AST-local:
\texttt{number\_replacer} (Num/Constant, CE partial, $\bigtriangleup$ literal values
only); \texttt{binary\_operator\_replacement} (BinOp);
\texttt{comparison\_operator\_replacement} (Compare);
\texttt{boolean\_replacer} (BoolOp); \texttt{unary\_operator\_replacement}
(UnaryOp); \texttt{keyword\_replacer} (NameConstant, CE keyword, $\bigtriangleup$
Bool literal); \texttt{break\_continue} (Break/Continue);
\texttt{exception\_replacer}; \texttt{remove\_decorator};
\texttt{variable\_inserter}; \texttt{variable\_replacer};
\texttt{no\_op}; \texttt{zero\_iteration\_for\_loop} (For). None correspond to OS / HP / TF / SI:

{%
\begin{longtable}[]{@{}
  >{\raggedright\arraybackslash}p{(\linewidth - 4\tabcolsep) * \real{0.3333}}
  >{\raggedright\arraybackslash}p{(\linewidth - 4\tabcolsep) * \real{0.3333}}
  >{\raggedright\arraybackslash}p{(\linewidth - 4\tabcolsep) * \real{0.3333}}@{}}
\caption{Per-class operator specialisations.}\label{tab:p2-18}\\
\toprule\noalign{}
\begin{minipage}[b]{\linewidth}\raggedright
Semantic class
\end{minipage} & \begin{minipage}[b]{\linewidth}\raggedright
Tool support
\end{minipage} & \begin{minipage}[b]{\linewidth}\raggedright
Coverage
\end{minipage} \\
\midrule\noalign{}
\endfirsthead
\endhead
\bottomrule\noalign{}
\endlastfoot
\textbf{OS} API replacement & None & $\bigtriangleup$ 88.33\% disjoint (incidental
low-complexity hits) \\
\textbf{HP} & None & $\times$ tool inexpressible (0/72) \\
\textbf{TF} numerical-method order & None & $\times$ tool inexpressible
(0/54) \\
\textbf{SI / CF} structural injection & None & $\times$ tool inexpressible
(0/33) \\
\end{longtable}
}

\textbf{Operator-level conclusion.} All 12 default classes remain
AST-local (BinOp / Compare / BoolOp / UnaryOp / NameConstant / Subscript
/ If / For / Break / Continue / decorator / except). Categorically, no
entry recognises sklearn / scipy hyperparameter semantics (HP),
numerical-method order (TF), or control-flow intent (SI / CF). For OS,
low-complexity sub-expressions can be incidentally hit by BinOp;
empirical 11.67\% (the main manuscript's AST-overlap audit) is consistent with 88.33\% AST-disjointness
lower bound.

\subsection{Syntactic-Mutation Pilot Study}\label{b.3-cosmic-ray-a1-single-put-pilot-pre-12-put-empirical}

Before the 12-PUT generalisation in the main manuscript's AST-overlap audit, a single-PUT lightweight
empirical used \path{scripts/run_cosmic_ray_a1.sh} on a1 (file
934 B, AST-bound). Outcome: mutants\_generated \textasciitilde tens; $\geq$
90\% belong to BinOp / Compare / 12 default classes; small
killed-by-\texttt{test\_a1.py} ratio (PUT unit tests are
output-shape/type sanity checks; most cosmic-ray mutants not killed), an
empirical manifestation of the semantic-mutation criteria in the main manuscript: syntactic-tool mutants
are not aligned with MR-violation detection. The 12-PUT generalisation
(the main manuscript's AST-overlap audit) supersedes this pilot but corroborates the same conclusion at
scale.

\subsection{Expected Root-Cause Risk Profile}\label{b.4-expected-lrca-risk-profile}

\textbf{PUT-class $\times$ LRCA-layer risk weights.} A numeric: C2 dominant
($\star\star\star$ on L1 tolerance). B probabilistic: C4 dominant ($\star\star\star$ on L3). C
surrogate: C3 dominant ($\star\star\star$ on L2 OOD; $\star\star$ tolerance). D ML: C3 + C4
mixed ($\star\star\star$ L2; $\star\star$ L3).

\textbf{Operator-PUT root-cause hotspots (expected).} A: mut\_C/G/F
dominantly C2 (numerical tolerance), mut\_M/T C1 (true semantic). B:
dominantly C4 (statistical-assumption violations). C: mut\_C/M
dominantly C2/C1 on c1/c2, mut\_T/F dominantly C3 on c1/c2; c3 (NN-Surr)
row dominantly C5 (artefact). D: d1 row dominantly C5 (training-noise
artefact); d2/d3 dominantly C1 / C3 mixed.

\textbf{Expected suspect\_share thresholds.} A: 0.10-0.20 (acceptance $\leq$
0.25); B: 0.20-0.35 ($\leq$ 0.40); C: 0.20-0.30 ($\leq$ 0.35); D: 0.25-0.40 ($\leq$
0.45). 60-cell average: $\leq$ 0.20 (= H4; acceptance $\leq$ 0.25).

\subsection{Engineering-significance mapping}\label{b.5-engineering-significance-mapping}

The aligned slice $k = k^\star(i)$ is the H2 threshold-test slice; the
cross slice $k \neq k^\star(i)$
is control; 9 vacant cells ($\circ$, no MR exercised) are not formally
adjudicated and are repurposed in the main manuscript's R\_sem/R\_kill decoupling discussion as descriptive evidence for
R\_sem / R\_kill decoupling. The v3 $\to$ v4 micro-change
($\Delta\delta = -0.009$, paired-role bootstrap 95\% CI [-0.238, 0.207])
attributes to LLM-source diversity under a fixed prompt; this is the
source-axis null contrast underpinning the main manuscript's aligned-versus-cross results.

\subsection{Default-Operator Overlap between Syntactic Mutation Tools}\label{b.6-mutmut-vs-cosmic-ray-default-operator-overlap}

Reviewer-requested manual operator-class cross-reference between
cosmic-ray's 13 default operators (cosmic-ray 8.4.6,
\path{cosmic_ray.operators} package) and mutmut's 14 default
operators (mutmut current \texttt{main},
\path{src/mutmut/mutation/mutators.py}). Both default-operator sets
are first-order AST-local replacements; neither implements a
cross-function-boundary, domain-knowledge-dependent, or
algorithmic-class-changing mutation (the three necessary conditions in the semantic-mutation criteria of the main manuscript
for semantic mutation).

{\tiny\def\LTcaptype{table} % do not increment counter
\begin{longtable}[]{@{}
  >{\raggedright\arraybackslash}p{(\linewidth - 8\tabcolsep) * \real{0.1800}}
  >{\raggedright\arraybackslash}p{(\linewidth - 8\tabcolsep) * \real{0.2200}}
  >{\raggedright\arraybackslash}p{(\linewidth - 8\tabcolsep) * \real{0.2200}}
  >{\raggedright\arraybackslash}p{(\linewidth - 8\tabcolsep) * \real{0.1600}}
  >{\raggedright\arraybackslash}p{(\linewidth - 8\tabcolsep) * \real{0.2200}}@{}}
\caption{Mutmut vs cosmic-ray default-operator overlap.}\label{tab:p2-19}\\
\toprule\noalign{}
\begin{minipage}[b]{\linewidth}\raggedright
Operator class
\end{minipage} & \begin{minipage}[b]{\linewidth}\raggedright
cosmic-ray default
\end{minipage} & \begin{minipage}[b]{\linewidth}\raggedright
mutmut default
\end{minipage} & \begin{minipage}[b]{\linewidth}\raggedright
Reaches the semantic-mutation criteria (a/b/c)?
\end{minipage} & \begin{minipage}[b]{\linewidth}\raggedright
Semantic-class reachability
\end{minipage} \\
\midrule\noalign{}
\endfirsthead
\endhead
\bottomrule\noalign{}
\endlastfoot
Numeric literal $\pm$1 & \texttt{number\_replacer} &
\texttt{operator\_number} & (a) no, (b) no, (c) no & partial CE only
(e.g.~\texttt{\_RHO=28.0\ $\to$\ 27.5}) \\
Binary arithmetic swap (+, $-$, $\times$, $\div$) &
\texttt{binary\_operator\_replacement} & \texttt{operator\_swap\_op}
(binary subset) & (a) no, (b) no, (c) no & partial OS only
(incidental) \\
Comparison swap (\textless, $\leq$, \textgreater, $\geq$, ==, !=) &
\texttt{comparison\_operator\_replacement} & \texttt{operator\_swap\_op}
(compare subset) & (a) no, (b) no, (c) no & none \\
Boolean swap (and $\leftrightarrow$ or) & \texttt{boolean\_replacer} &
\texttt{operator\_swap\_op} (bool subset) & (a) no, (b) no, (c) no &
none \\
Unary remove / swap & \texttt{unary\_operator\_replacement} &
\texttt{operator\_remove\_unary\_ops}, \texttt{operator\_swap\_op} & (a)
no, (b) no, (c) no & none \\
Keyword swap (True/False, etc.) & \texttt{keyword\_replacer} &
\texttt{operator\_keywords}, \texttt{operator\_name} & (a) no, (b) no,
(c) no & none \\
break $\leftrightarrow$ continue / return & \texttt{break\_continue} &
\texttt{operator\_keywords} (subset) & (a) no, (b) no, (c) no & none \\
Exception class swap & \texttt{exception\_replacer} & (none) & (a) no,
(b) no, (c) no & none \\
Decorator removal & \texttt{remove\_decorator} & (none) & (a) no, (b)
no, (c) no & none \\
Variable-binding insert / replace & \texttt{variable\_inserter},
\texttt{variable\_replacer} & \texttt{operator\_arg\_removal},
\texttt{operator\_assignment} & (a) no, (b) no, (c) no & none \\
no\_op insertion & \texttt{no\_op} & (none) & (a) no, (b) no, (c) no &
none \\
Zero-iteration \texttt{for} loop & \texttt{zero\_iteration\_for\_loop} &
(none) & (a) no, (b) no, (c) no & none \\
String content tweak (XX prefix, case flip) & (none) &
\texttt{operator\_string} & (a) no, (b) no, (c) no & none \\
Lambda body $\to$ \texttt{None}/\texttt{0} & (none) &
\texttt{operator\_lambda} & (a) no, (b) no, (c) no & none \\
Dict kwarg name (XX prefix) & (none) &
\texttt{operator\_dict\_arguments} & (a) no, (b) no, (c) no & none \\
Sym./asym. string-method swap & (none) &
string-method symmetric swap; string-method asymmetric swap & (a) no,
(b) no, (c) no & none \\
Augmented $\to$ simple assignment & (none) &
\texttt{operator\_augmented\_assignment} & (a) no, (b) no, (c) no &
none \\
\texttt{match} case removal & (none) & \texttt{operator\_match} & (a)
no, (b) no, (c) no & none \\
\end{longtable}
}

\textbf{Aggregate.} Cosmic-ray $\cap$ mutmut: \textasciitilde9 of 13
cosmic-ray operators have a near-equivalent in mutmut (numeric, binary,
compare, boolean, unary, keyword, break-continue, variable, plus partial
overlap on string-prefix vs \texttt{XX}); 4 cosmic-ray operators are
unique (exception, decorator, no\_op, zero-iteration-for); 6 mutmut
operators are unique (string content, lambda, dict-kwarg, string-method
swap, aug-assignment, match-case). The union is approximately 21
distinct first-order AST-local operator classes. None of the 21 can:

\begin{itemize}

\item
  \begin{enumerate}
  
  \item
    Cross a function-call or module-import boundary (e.g.~replace
    \texttt{det(M)} with \texttt{sum(np.diag(M))} requires knowing the
    equivalence on diagonal matrices, outside any default-operator
    scope);
  \end{enumerate}
\item
  \begin{enumerate}
  \setcounter{enumi}{1}
  
  \item
    Depend on domain knowledge for legality (e.g.~a hyperparameter swap
    that knows \texttt{noise\_level=1e-4} is the load-bearing knob in a
    Gaussian Process Regression PUT requires PUT-class awareness);
  \end{enumerate}
\item
  \begin{enumerate}
  \setcounter{enumi}{2}
  
  \item
    Change the algorithmic class (e.g.~swap
    \texttt{polynomial\ $\to$\ spline\ basis} rewrites the surrogate's
    mathematical structure).
  \end{enumerate}
\end{itemize}

The mutmut / cosmic-ray null intersection on the main manuscript's AST-overlap audit's HP/SI/TF classes
(zero overlap) therefore reflects a \textbf{structural property of
first-order AST-local mutation} rather than a tool-specific limitation,
supporting the preventive-defence framing in the main manuscript.

% [thematic break removed]

\section{Experimental Procedure Details}\label{c.-experimental-procedure-details}

\subsection{Cross-Source Generation Protocol}\label{c.1-cross-source-v4-protocol-full-specification}

\textbf{Two-stage ablation.}

{%
\begin{longtable}[]{@{}
  >{\raggedright\arraybackslash}p{(\linewidth - 6\tabcolsep) * \real{0.2500}}
  >{\raggedright\arraybackslash}p{(\linewidth - 6\tabcolsep) * \real{0.2500}}
  >{\raggedright\arraybackslash}p{(\linewidth - 6\tabcolsep) * \real{0.2500}}
  >{\raggedright\arraybackslash}p{(\linewidth - 6\tabcolsep) * \real{0.2500}}@{}}
\caption{Cross-source generation protocol.}\label{tab:p2-20}\\
\toprule\noalign{}
\begin{minipage}[b]{\linewidth}\raggedright
Version
\end{minipage} & \begin{minipage}[b]{\linewidth}\raggedright
Mutant pool
\end{minipage} & \begin{minipage}[b]{\linewidth}\raggedright
c-class primary MR
\end{minipage} & \begin{minipage}[b]{\linewidth}\raggedright
Use
\end{minipage} \\
\midrule\noalign{}
\endfirsthead
\endhead
\bottomrule\noalign{}
\endlastfoot
v3 & Single-source (Claude Opus 4.6) & P-ord. (pre-registered) &
H2 baseline \\
v4 & \textbf{Cross-source} (Claude Opus 4.6 + GPT-5.4 + DeepSeek chat) &
P-ord. (held at pre-registered) & Isolate LLM-source diversity \\
\end{longtable}
}

\textbf{Protocol} (\path{scripts/cross_source_campaign.py}). (a) For
each (PUT, operator), Claude / GPT / DeepSeek run K = 3 trials with
identical prompt (the cross-source protocol in the main manuscript), temperature 0.7; \texttt{source\_tag}
propagated to filenames. (b) \textbf{V1-V4 mechanical validation}
(\path{src/p2/mutators/validation.py}): V1 syntax
(\texttt{ast.parse}), V2 executability, V3 non-triviality
(\texttt{\textbar{}y\_mutant\ -\ y\_original\textbar{}\ \textgreater{}\ 1e-6}),
V4 signature consistency. (c) \textbf{DeepSeek model:}
\texttt{deepseek-chat} (113 tokens/call, 1.8 s), not
\texttt{deepseek-v4-pro} (340 tokens/call, 11 s), quality-equivalent
on a2\_OS1 dry-run. (d) \textbf{Pool capacity:} 37 ops $\times$ 3 sources $\times$ 3
trials = 333 attempts; V1-V4 pass rate 89.5\%; 298 confirmed mutants, of
which 292 enter the v4 pool after final-stage rejection of duplicates and
QA failures. Sources contribute nearly equally: Claude 101 / GPT 98 / DeepSeek 99
(Phase A engineering finding). (e) \textbf{Sampling}
(\path{scripts/build_pools.py}, POOL\_VERSION = v4): per-PUT 30 max,
measured mean 24.3, range 10-30. c1 (GPR) has only 10 because c1\_HP1
/ c1\_CE1 have V1-V4 near-zero pass (WhiteKernel noise\_level 1e-4 $\to$
1e-1 perturbation has minimal output impact, triggers V3 non-trivial
failure, which supports the main manuscript's R\_sem/R\_kill decoupling discussion).

\textbf{Protocol asymmetry.} v4 does not invoke reviewer LLM (cost
/ speed priority; a follow-up study will rerun dual-blind on the v4
grid, \textasciitilde\$5-8 USD); v3 used the original Phase-1 dual-blind
(Claude gen + GPT-5.4 review + DeepSeek arbitration). A fraction of the
v3 $\to$ v4 source-axis shift may be slight v4 quality decline rather than
LLM-source diversity.

\subsubsection{Differential prompt protocol}\label{c.1.1-differential-prompt-protocol-future-work-commitment}

The v3 $\to$ v4 micro-change of -0.009 (paired-role bootstrap 95\% CI
[-0.238, 0.207]) may reflect ``three LLMs under
identical prompt'' rather than ``true upper bound of LLM diversity''.
Separation requires \emph{one tailored differential prompt per LLM}
(skeleton: \path{scripts/run_differential_prompt.py}).

\textbf{Design (2$\times$3 factorial, within-(PUT, operator)).} Factor A:
prompt template, 3 levels: \textbf{V\_canonical} (control, identical to
v4), \textbf{V\_persona} (per-LLM identity: Claude ``numerical-analysis
editor''; GPT ``scientific-software refactorer''; DeepSeek ``library-API
substitution specialist''), \textbf{V\_cot} (Claude
\texttt{\textless{}thinking\textgreater{}} tags; GPT ``step by step'';
DeepSeek stepped-reasoning). Factor B, LLM source (Claude / GPT-5.4 /
DeepSeek chat). K = 3 per cell; total 37 $\times$ 3 $\times$ 3 $\times$ 3 = \textbf{999
trials}.

\textbf{Exit criteria.} If max(delta) $-$ min(delta) \textless{} 0.05: the
v3 $\to$ v4 $-$0.009 source-axis null is robust. If $\geq$ 0.05 with prompt variant
pushing delta past 0.474: H2 revised from ``not met'' to
``prompt-conditional met''. If $\geq$ 0.05 but all delta \textless{} 0.474:
prompt sensitivity exists but H2 unchanged.

\textbf{Resources.} API \textasciitilde\$18-30 USD; wall 3-4 h at
concurrency 4; V1-V4 only, no reviewer LLM.

\subsection{Manual Pilot and Initial Prescreen}\label{c.2-manual-pilot-lineage-llm-client-configuration-and-l0-prescreen}

The original cross-source protocol used 60\% multi-LLM consensus (Claude Opus
4.6 + GPT-5.4 + DeepSeek chat unanimous) plus 40\% manual injection by
scientific-computing researchers. Each (mut\_j, S\_i) pair carried one
prompt template containing PUT source, semantic intent of mut\_j,
prohibition against revealing specific MRs (avoids prompt leakage), and
output requirements (syntactically correct, executable, single-point
modification, diff \textless{} 10 lines). Reproducibility parameters:
temperature 0.3 in manual-pilot stage (v4 cross-source uses 0.7 per
§C.1); fixed seed; 5 candidates per prompt with 2-3 retained.

\textbf{Double-blind review (Scheme C).} Generator LLM-G (Claude Opus)
and reviewer LLM-R (GPT-5.4) are cross-source with strict role
separation. LLM-R sees only PUT original + mutant code, unaware of
mut\_j category / MR content / generator identity, outputting a triple
(syntactic, executable, semantic-fault-injected). Double-confirmed
mutants enter pool; inconsistencies enter manual arbitration ($\leq$ 10\%).
From the double-confirmed pool, 20\% is sampled for manual review;
manual-vs-LLM-R inconsistency \textgreater{} 10\% triggers full manual
downgrade. Same-source bias mitigation stratifies LLM-R by PUT class.
Prompt-injection mitigation disables RAG / web search on the API path.

\textbf{L0 pool prescreen.} Each mutant passes static syntax (linters /
type checkers), a unit self-test (simple inputs produce finite output),
and double-blind sign-off. Target: 30-50 candidates per cell $\to$ 10-15
retained after prescreen.

\subsection{Three-Layer Root-Cause Diagnosis}\label{c.4-lrca-three-layer-execution-full}

Decision tree, multi-label priority (C5 \textgreater{} C4 \textgreater{}
C3 \textgreater{} C2 \textgreater{} C1), and output (\texttt{C1\_share},
\texttt{suspect\_share}) are stated in §A.2; the full pseudocode is
reproduced there. LRCA does \textbf{not} modify SMS; killed set is not
filtered. The main manuscript's LRCA-mass results use the best calibrated combination
(\texttt{ood\_band\ =\ 0.02,\ tolerance\_multiplier\ =\ 3.0,\ repeats\ =\ 20});
default-threshold results retained in \texttt{lrca\_60cell\_v3.json} as
control. Interface with the preventive-defence risk profile in the main manuscript: L0 prescan $\leftrightarrow$ §C.2 dual-LLM
double-blind; L1 tolerance $\leftrightarrow$ N = 20 repetition subprocess; L2 OOD $\leftrightarrow$
input-distribution definition; L3 statistical baseline $\leftrightarrow$ pre-check
protocol.

\subsubsection{Pilot calibration}\label{c.4.1-pilot-calibration-37-operators-k-1020}

A 2026 Q3 operator-level pilot ran 12 PUTs $\times$ 37 operators (36
primary PUT-class operator pairs plus one CF structural-injection
specialisation; the conceptual semantic family count remains five; 12
\texttt{is\_key=True} at K=20, 25 at K=10; 470 trials total). Pilot
precursor quantities: R\_sem (V1-V4 $\wedge$ operator\_match pass rate),
D\_impl (median pairwise AST + literal + identifier Jaccard distance
among confirmed mutants), R\_kill (proportion of confirmed mutants
killed by AVP under that PUT's primary MR). Aggregate (N=37): R\_sem
median 0.50 mean 0.468 (0/37 with R\_sem=0); D\_impl median 0.42 mean
0.392 (1/37 with D\_impl $\approx$ 0); R\_kill median 0.00 mean 0.189. All 37
operators produced $\geq$ 1 V1-V4 passing mutant, a necessary
precondition for H1; the pre-registered H1 threshold itself ($\geq$ 5
non-equivalent mutants on $\geq$ 9/12 PUTs for $\geq$ 4 of 5 operator
families) is nonetheless not met, as reported in the main text.

\textbf{Key R\_sem / R\_kill decoupling pattern.} R\_sem $\geq$ 0.9 $\wedge$ R\_kill
= 0 combinations all on HP / TF / OS operators in c / d classes (e.g.,
c1\_CE1 noise 1e-4$\to$1e-1 with R\_sem 1.00, R\_kill 0.00; c3\_HP1
relu$\to$tanh, c3\_TF1 max\_iter 1000$\to$5, d1\_HP1 MLP $\alpha$, d3\_HP1 LR C).
R\_sem-high $\wedge$ R\_kill = 1 combinations all on CE / OS operators in a / b
classes (a2\_CE1 LU det, a2\_OS1 prod$\to$sum, b1\_OS1 $\alpha$/$\beta$ swap, d1\_TF1
label flip). This is empirical evidence for R\_sem / R\_kill decoupling at cell level.

% [thematic break removed]

\section{Statistical Analysis Details}\label{d.-statistical-analysis-details}

\subsection{Reporting Conventions and Primary-Table Construction}\label{d.1-sms-variants-and-construction-lemmas}

An earlier draft defined a separate \texttt{SMS\_unfiltered} variant
``without LRCA filtering''. Under the final definitions LRCA never
filters the killed set, so that variant coincides with SMS identically
and has been removed; the diagnostic secondary view is the C1-weighted
$\mathrm{SMS}_{\mathrm{C1}}$ of the main text, which is explicitly a
product of two reported quantities rather than a second mutation score.

\textbf{Primary-table construction.}

{%
\begin{longtable}[]{@{}
  >{\raggedright\arraybackslash}p{(\linewidth - 4\tabcolsep) * \real{0.3333}}
  >{\raggedright\arraybackslash}p{(\linewidth - 4\tabcolsep) * \real{0.3333}}
  >{\raggedright\arraybackslash}p{(\linewidth - 4\tabcolsep) * \real{0.3333}}@{}}
\caption{SMS variants and construction lemmas.}\label{tab:p2-21}\\
\toprule\noalign{}
\begin{minipage}[b]{\linewidth}\raggedright
RQ
\end{minipage} & \begin{minipage}[b]{\linewidth}\raggedright
Primary statistics
\end{minipage} & \begin{minipage}[b]{\linewidth}\raggedright
Reporting format
\end{minipage} \\
\midrule\noalign{}
\endfirsthead
\endhead
\bottomrule\noalign{}
\endlastfoot
RQ3 & inst\_rate, equiv\_rate, C1\_share, survive\_rate & 60-cell
heatmap + 4-class marginals \\
RQ4 & aligned-SMS vs cross-SMS; sparse $\circ$ vs dense \textbullet\textbullet equiv\_rate &
Cliff's delta + odds ratio + 95\% bootstrap CI \\
RQ5 & $\Delta$SMS\_c (c $\in$ \{A,B,C,D\}); CV($\Delta$SMS) & sign test (df=3) +
descriptive forest plot \\
RQ5 (desc.) & Spearman $\rho$ + Kendall $\tau$ for SMS vs PC & scatter + dual-metric
ranking comparison \\
\end{longtable}
}

\textbf{Multiple comparisons.} No cell-level family of p-values is
claimed across the 60 cells (main text, statistical pipeline); the
per-class Friedman tests carry a Bonferroni $\times$ 4 correction.
N = 20 bootstrap intervals:
1,000-iteration 95\% CI (B = 10,000 for the headline H2 delta CI).

\subsection{Effect-Size Sensitivity and Transformation Invariance}\label{d.2-cliffs-delta-cutoff-sensitivity-and-logit-transformation-invariance}

\textbf{H4 cutoff sensitivity.} Dense grid scan (cutoff
$\in$ \{0.05, 0.10, \ldots, 0.50\}, step 0.05) on v4 data
(\path{scripts/h5_sensitivity.py},
\path{data/results/h5_sensitivity_v4.json}), re-based on the 15
evaluable cells per the main text's \texttt{NA} convention (zero-kill
cells carry no defined suspect\_share):

{%
\begin{longtable}[]{@{}lll@{}}
\caption{Cutoff sensitivity of the per-cell suspect-share pass count
over evaluable cells.}\label{tab:p2-22}\\
\toprule\noalign{}
cutoff & cells at or below cutoff & share of evaluable \\
\midrule\noalign{}
\endfirsthead
\endhead
\bottomrule\noalign{}
\endlastfoot
0.05 & 12 / 15 & 80.0\% \\
0.10 & 12 / 15 & 80.0\% \\
0.15 & 12 / 15 & 80.0\% \\
\textbf{0.20 (paper)} & \textbf{12 / 15} & \textbf{80.0\%} \\
0.25 & 12 / 15 & 80.0\% \\
0.30 & 12 / 15 & 80.0\% \\
0.35 & 12 / 15 & 80.0\% \\
0.40 & 12 / 15 & 80.0\% \\
0.45 & 13 / 15 & 86.7\% \\
0.50 & 13 / 15 & 86.7\% \\
\end{longtable}
}

\textbf{Conclusion.} The per-cell pass count is flat across cutoffs
$\in$ {[}0.05, 0.40{]}: the evaluable cells' suspect\_share
distribution is strongly bimodal (twelve cells at or near 0, a small
high-suspect tail carrying most of the pooled suspect mass), so the
specific cutoff 0.20 is \textbf{not load-bearing}. An earlier draft
reported this sweep over all 60 cells (12/60 flat) by coding zero-kill
cells as suspect\_share = 1; that coding is withdrawn with the H4
estimand discussion in the main text, and no pass/fail verdict attaches
to this table.

\textbf{Logit-transformation invariance.} Cliff's delta is rank-based
(function of $U / (n_1 \cdot n_2)$), mathematically invariant under any strictly
monotone transformation (Romano 2006). Logit is strictly monotone on (0,
1), so $\delta_{\mathrm{logit}} \equiv \delta_{\mathrm{raw}}$ is a construction result. \textbf{Not
additional robustness evidence}; H2 robustness threats come from the main manuscript's 60-cell distribution results
zero-mass dominance, not metric-scale choice.

\textbf{Vacant-cell sensitivity.} Excluding the 9 vacant cells (cells
with no exercised MR) leaves the aligned-versus-cross verdict unchanged.
In v4 the analysis becomes n = 11 aligned and n = 40 cross with
$\delta = 0.323$ (PUT-cluster 95\% CI [0.045, 0.599]). The exclusion
does not move H2 across the pre-registered large-effect threshold.
For the complete primary grid, 100,000 PUT-cluster resamples (seed 42)
give \(\delta=0.314\), 95\% CI [0.045, 0.594], with 0.00964 of
resamples at or below zero
(\path{scripts/compute_rq2_cluster_bootstrap.py};
\path{data/results/rq2_cluster_bootstrap_v4.json}).

\textbf{Primary-MR withdrawal permutation-null.} The auxiliary c-class
check pools 15 SMS values, shuffles labels over the $3 \times 5$ (PUT,
MR) grid, and applies the max-over-five withdrawal rule with 10,000
permutations (seed 42). The observed c-class mean is 0.314 versus null
mean 0.349, one-sided $p = 0.989$
(\path{data/results/c_class_permutation_v4.json}). This supports treating
the primary-MR convention as a protocol choice rather than as evidence of
a hidden c-class inflation effect.

\subsection{Power Analysis under Observed and Stipulated Effects}\label{d.3-power-analysis-plug-in-and-stipulated-alternative}

\textbf{Plug-in bootstrap.} With-replacement sample from observed
aligned (n=12) and cross (n=48) v4 SMS pools; N\_sim = 5,000, seed = 42
(\path{scripts/compute_rq2_power.py}, \path{rq2_power_v4.json}):

{%
\begin{longtable}[]{@{}lll@{}}
\caption{Observed-distribution exceedance probabilities.}\label{tab:p2-23}\\
\toprule\noalign{}
Threshold & Interpretation & Exceedance probability \\
\midrule\noalign{}
\endfirsthead
\endhead
\bottomrule\noalign{}
\endlastfoot
delta \textgreater{} 0.000 & Any-effect & \textbf{0.978} \\
delta \textgreater{} 0.147 & Small & 0.849 \\
delta \textgreater{} 0.330 & Medium & 0.457 \\
\textbf{delta \textgreater{} 0.474} & \textbf{Large (H2)} &
\textbf{0.155} \\
\end{longtable}
}

Sample-size sweep ($n_{\mathrm{aligned}} \in \{6, 12, \ldots, 60\}$, $n_{\mathrm{cross}} = 4 \times n_{\mathrm{aligned}}$):
the probability for delta \textgreater{} 0 is 0.909 at n=6, 0.978 at
12, and 0.994 at 18. These are observed-distribution exceedance
probabilities, not power against a stipulated alternative. Observed
$\delta_{\mathrm{v3}} = 0.323 < 0.474$ and observed
$\delta_{\mathrm{v4}} = 0.314 < 0.474$; increasing sample only narrows CI.

\textbf{Stipulated-alternative.} Implemented in
mixture-weight (\path{scripts/compute_rq2_power_stipulated.py},
N\_sim = 2,000): SMS has heavy ties at zero (45/60 = 0), raw shifting is
discontinuous (any $\varepsilon$ \textgreater{} 0 jumps delta from 0.314 to 0.74).
Mixture: aligned' = (probability w) sample from (observed\_aligned +
0.001) + (probability 1$-$w) sample from observed\_aligned, calibrate w so
E{[}delta{]} $\approx$ 0.474. Calibrated w = 0.373, realised E{[}delta{]} =
0.470.

{%
\begin{longtable}[]{@{}lll@{}}
\caption{Power analysis under observed and stipulated effects.}\label{tab:p2-24}\\
\toprule\noalign{}
Stipulated truth & Criterion & Power \\
\midrule\noalign{}
\endfirsthead
\endhead
\bottomrule\noalign{}
\endlastfoot
0.470 (target 0.474) & $\hat{\delta} \geq 0.474$ (point-estimate) & \textbf{0.504} \\
0.470 (target 0.474) & 95\% CI lower \textgreater{} 0 (any-effect) & 0.895 \\
\end{longtable}
}

At the nearest attainable mixture, \(E[\delta]\approx0.470\), the
(12, 48) design returns $\hat{\delta} \geq
0.474$ in \textasciitilde50\% of replications. This
\emph{supports} the main manuscript's aligned-versus-cross framing that ``not met'' is point-estimate, not
effect-size.

\subsection{Per-Class Friedman Tests}\label{d.4-friedman-per-class-discussion-of-small-n-concordance}

Per-class Friedman (3 SUTs $\times$ 5 MRs, v4 primary pool) with Bonferroni
$\times$ 4: a chi\^{}2=4.00
raw 0.406 adj 1.000, W=0.333; b chi\^{}2=10.34 raw \textbf{0.035} adj
\textbf{0.140}, W=\textbf{0.862}; c chi\^{}2=5.60 raw 0.231 adj 0.924,
W=0.467; d chi\^{}2=5.00 raw 0.287 adj 1.000, W=0.417. After correction,
no per-class result remains significant at family-wise alpha = 0.05.
Even b-class W = 0.862 (Cohen 1988 large concordance) is insufficient at
N = 3 per class. \textbf{H3 verdict rests on the main manuscript's cross-class consistency sign test}; per-class
Friedman is sensitivity / descriptive only. Friedman main effect on the full
60 cells (v4: chi\^{}2 = 16.76, p = 0.0022; v3: 15.30, 0.0041) confirms
MR rank differences but
is logically independent from H3 cross-class consistency.

\textbf{Mixed-effects unavailability.} Primary
\texttt{sms\ \textasciitilde{}\ C(class)\ +\ C(operator)\ +\ C(class):C(operator)\ +\ (1\ \textbar{}\ put)}
returned Singular (insufficient column rank for class $\times$ operator
interaction; N=60 too small for 11-d fixed-effects + 12 PUT random
intercepts). Fallback
\texttt{sms\ \textasciitilde{}\ C(class)\ +\ C(operator)\ +\ (1\ \textbar{}\ put)}
had PUT random-intercept variance hit boundary 0 (degenerates to OLS).
Fallback class p-values (descriptive only): b vs a 0.275 / c vs a 0.892
/ d vs a 0.991. RQ5 conclusion shifts to four-piece presentation (class
means + sign test + Friedman + forest plot).

\subsection{Pattern-coverage operationalisation}\label{d.5-pattern-coverage-operationalisation}

Per SUT, (MR\_k, R\_outcome $\in$ \{True, False\}) binary-tuple coverage: 5
MRs $\times$ 2 outcomes = 10 cells, PC = \#triggered / 10. Simplest baseline
(per-SUT granularity, not distinguishing mutants). PC range over 12
SUTs (v4): {[}0.500, 1.000{]}, mean 0.750
(\texttt{rq4\_pattern\_coverage\_v4.json}).
Pairing per PUT: Spearman rho = 0.163; Kendall tau = 0.136
(descriptive).

\textbf{Power caveat.} At n = 12, Spearman 95\% CI $\approx$ {[}-0.5, +0.6{]}
(Fisher-z approximation); cannot distinguish ``zero correlation'' from ``moderate positive'' or
``moderate negative''. The descriptive reading is ``no correlation detected
at n = 12'', not ``correlation does not exist''.

\textbf{Within-class descriptive.} In the b class, b2 (PC = 1.0) has
mean SMS 0.133 while b1 (PC = 0.7) has 0.20: higher PC, lower SMS. In
the c class the ordering reverses: c3 (PC = 1.0) has mean SMS 0.148,
c1 (PC = 0.9) has 0.12, and c2 (PC = 0.8) has 0.0. The within-class
PC-SMS ordering is therefore not stable across classes, consistent with
the headline finding of no detectable correlation at n = 12; an earlier
draft reported v3 values here and read the b-class pattern as
uniformly ``higher PC, lower SMS'', which the v4 numbers do not
support. A confirmatory study would need n $\geq$ 30 and a PC
definition that incorporates the mutant dimension.

% [thematic break removed]

\section{Stakeholder Deployment Considerations}\label{e.-stakeholder-deployment-considerations-single-output-kernels-scope}

\subsection{Air-Gap Deployment Constraints}\label{e.1-air-gap-incompatibility-full-justification}

The stakeholder-analysis workflow in the main manuscript depends on external LLM API calls (Claude / GPT /
DeepSeek), incompatible with most regulated air-gapped V\&V workflows:

{%
\begin{longtable}[]{@{}
  >{\raggedright\arraybackslash}p{(\linewidth - 4\tabcolsep) * \real{0.3333}}
  >{\raggedright\arraybackslash}p{(\linewidth - 4\tabcolsep) * \real{0.3333}}
  >{\raggedright\arraybackslash}p{(\linewidth - 4\tabcolsep) * \real{0.3333}}@{}}
\caption{Air-gap deployment constraints.}\label{tab:p2-25}\\
\toprule\noalign{}
\begin{minipage}[b]{\linewidth}\raggedright
Domain
\end{minipage} & \begin{minipage}[b]{\linewidth}\raggedright
Standard
\end{minipage} & \begin{minipage}[b]{\linewidth}\raggedright
Air-gap requirement
\end{minipage} \\
\midrule\noalign{}
\endfirsthead
\endhead
\bottomrule\noalign{}
\endlastfoot
Nuclear & IEC 60880 & Air-gapped build for safety-critical software \\
Aerospace & DO-178C & Tool-qualified offline build for DAL-A/B; LLM API
outside qualified-tool envelope \\
Medical & IEC 62304 & Class C software lifecycle requires offline
reproducible build \\
Automotive & ISO 26262 & ASIL-D requires offline tool chain; cloud LLM
non-compliant \\
\end{longtable}
}

\textbf{Mitigation paths.} (a) Self-hosted open-weight LLMs (Llama
/ DeepSeek local inference), capability gap vs API-served frontier
models; quantitative impact on SMS untested. (b) Offline-cached
pre-generated mutant pools with commit-hash + raw-response signature
locking, adopters reproduce a published pool offline (this paper's
package supports this), but generating new pools requires external LLM
access. SMS for air-gapped industrial deployment is a future research
direction.

\subsection{Quarterly batch audit workflow}\label{e.2-quarterly-batch-audit-workflow}

The original per-PR GitHub Actions YAML was removed because its PR-CI
threshold 0.10 contradicted the Appendix~E.3 audit threshold 0.20 and propagated
stakeholder-facing CI assumptions into adopters' pipelines.

\textbf{Replacement: quarterly batch audit:}

{%
\begin{longtable}[]{@{}
  >{\raggedright\arraybackslash}p{(\linewidth - 4\tabcolsep) * \real{0.3333}}
  >{\raggedright\arraybackslash}p{(\linewidth - 4\tabcolsep) * \real{0.3333}}
  >{\raggedright\arraybackslash}p{(\linewidth - 4\tabcolsep) * \real{0.3333}}@{}}
\caption{Quarterly batch audit workflow.}\label{tab:p2-26}\\
\toprule\noalign{}
\begin{minipage}[b]{\linewidth}\raggedright
Step
\end{minipage} & \begin{minipage}[b]{\linewidth}\raggedright
Description
\end{minipage} & \begin{minipage}[b]{\linewidth}\raggedright
Cost
\end{minipage} \\
\midrule\noalign{}
\endfirsthead
\endhead
\bottomrule\noalign{}
\endlastfoot
1 & Offline mutant-pool generation (per PUT 24-30 $\times$ 12 PUTs $\approx$ 300) & LLM
API \$5-15; \textasciitilde30 min \\
2 & \texttt{sms\_campaign.py\ -\/-track\ 2} & 4-core laptop
\textasciitilde10-20 min \\
3 & \texttt{run\_lrca.py} & \textless{} 1 min \\
4 & MR-design team review of MRs with aligned-cell SMS significantly
below the historical aligned mean (0.213) & 1-2 h meeting \\
\end{longtable}
}

Quarterly: API \textasciitilde\$10 + automation 30 min + manual 1-2 h $\approx$
0.5 person-day. Affordable. Per-PR gating not recommended (LLM latency
5-30 s + cost + non-determinism + air-gap incompatibility); adopters
needing per-PR should use coverage / unit-test gates. SMS does not
replace domain-expert MR physical-reasonableness judgement nor
product-requirements review.

\subsection{Verification-and-Validation Documentation}\label{e.3-vv-documentation-conceptual-complementarity-with-asme-vv-20-2009}

This appendix subsection makes \textbf{no normative claim} toward NRC,
FDA, or ISO 26262 review teams. Within single-output kernels scope, no
traceable mapping exists to current IEC 60880 / ISO 26262 / DO-178C /
ASME V\&V 20-2009 normative bodies.

V\&V documentation for scientific computing software (per ASME V\&V
20-2009 §3 and similar guides) requires quantifiable test-adequacy
evidence; current documentation often relies on code coverage + MR lists
+ SME signatures. SMS may appear as \textbf{research-grade supplementary
evidence} alongside coverage and MR lists: (a) aligned-cell SMS per
critical PUT; (b) LRCA three-layer diagnosis (C1 readout / C2-C5
attribution); (c) 60-cell matrix visualisation. Reviewers can
independently run \texttt{REPRODUCIBILITY.md}.

Original threshold recommendations (aligned-cell SMS $\geq$ 0.20 / 0.30 +
C1\_share $\leq$ 0.20) were removed: no normative backing in IEC / ISO /
ASME; misreads as enforce-ready. \textbf{SMS reported as descriptive
supplementary evidence}, with adequacy judgements left to V\&V reviewers
case-by-case.

Conceptually complementary to ASME V\&V 20-2009 §3 code verification
(numerical-solver correctness) vs SMS (MR-set fault-detection adequacy).
Substantial scale gap to multi-module CFD codes; incorporation into V\&V
standards body would require large-scale empirics + multi-year
dialogue with ASME V\&V 20 / IEEE 1012 / IEC 60880 committees.
\textbf{No advocacy that SMS enter any normative certification system in
2027.} SMS does not replace SME signatures, FMECA, PIRT, system-level
V\&V, or ASME V\&V 20 §3 numerical-solver verification, it is one
link in the evidence chain, not a single-point qualification.

% [thematic break removed]

\section{Threats to Validity: Detailed Mitigation}\label{f.-threats-to-validity-detailed-mitigation}

\subsection{Internal threats}\label{f.1-internal-threats}

\textbf{Hoeffding-style false-equivalence bound (referenced by
main-text Limitation 1).} For a mutant whose true per-sample
disagreement probability against the original program under $\mathcal{D}_P$ is
at least $p_0$ (at the E2 tolerance), the probability that all
$K_{\mathrm{eq}}$ independent E2 samples agree, producing a false
equivalence verdict, is at most $(1-p_0)^{K_{\mathrm{eq}}} \leq
\exp(-K_{\mathrm{eq}}\,p_0)$. At $K_{\mathrm{eq}} = 1000$: $p_0 =
0.005$ gives $\leq e^{-5} \approx 0.7\%$; $p_0 = 0.01$ gives $\leq
e^{-10} \approx 4.5 \times 10^{-5}$. The bound conditions on $p_0$ and
says nothing about mutants whose disagreement is concentrated below the
E2 tolerance; those are governed by the E1 arm and remain the residual
risk the K\_eq sweep (deferred) would quantify.

{%
\begin{longtable}[]{@{}
  >{\raggedright\arraybackslash}p{(\linewidth - 6\tabcolsep) * \real{0.2500}}
  >{\raggedright\arraybackslash}p{(\linewidth - 6\tabcolsep) * \real{0.2500}}
  >{\raggedright\arraybackslash}p{(\linewidth - 6\tabcolsep) * \real{0.2500}}
  >{\raggedright\arraybackslash}p{(\linewidth - 6\tabcolsep) * \real{0.2500}}@{}}
\caption{Internal threats.}\label{tab:p2-27}\\
\toprule\noalign{}
\begin{minipage}[b]{\linewidth}\raggedright
Threat area
\end{minipage} & \begin{minipage}[b]{\linewidth}\raggedright
Threat
\end{minipage} & \begin{minipage}[b]{\linewidth}\raggedright
Mitigation
\end{minipage} & \begin{minipage}[b]{\linewidth}\raggedright
Residual
\end{minipage} \\
\midrule\noalign{}
\endfirsthead
\endhead
\bottomrule\noalign{}
\endlastfoot
LLM generation & LLM generation reproducibility & §C.2 reproducibility parameters in
package; dual-LLM cross-source + 20\% manual; reproducibility rests on
archived raw responses and frozen mutant pools (no user-facing seed) & LLM training data may be updated
post-submission \\
Equivalence approximation & Probabilistic equiv (K\_eq=1000) & the main manuscript's equivalence-judgement section declares probabilistic
approximation; Hoeffding-style false-equiv bound & K\_eq sweep $\in$
\{500,1000,2000\} deferred to future work \\
Shared infrastructure & Shared-infrastructure dependency & AVP version pinned to the upstream commit hash;
package embeds AVP source; the arXiv infrastructure technical report is disclosed in the main text &
None within scope \\
Root-cause diagnosis & LRCA multi-label boundary & §A.2 priority
C5 \textgreater{} C4 \textgreater{} C3 \textgreater{} C2 \textgreater{} C1;
multi-label co-occurrence table in package; root\_cause is ``likely root
cause'' not definitive & C2-C5 may multi-trigger on B/D classes \\
Pool size & Mutant-pool size & Pool expanded to mean 17.4: delta moved 0.321 $\to$
0.323, CI narrowed; \textbf{effect size is intrinsic ceiling, not pool
dilution} & Some PUTs \textless{} 30 (cache-bound) \\
LLM non-determinism & LLM non-determinism & Multi-turn de-dup; K=10/20 repetitions; raw
prompt+response committed for direct reuse & Claude Opus subscription
lacks seed control \\
\textbf{Protocol asymmetry} & \textbf{Protocol-implementation gap v3 vs v4} & v4 used 3
providers (Claude 101 / GPT 98 / DeepSeek 99) without dual-blind; v4
LRCA-evaluable cells increase 12 $\to$ 15 and macro C1 share changes
0.821 $\to$ 0.837 & Complete separation requires a follow-up dual-blind rerun on
the v4 grid \\
\end{longtable}
}

The protocol-asymmetry threat (dual-blind vs V1-V4 only) applies to the
v3 $\to$ v4 source-axis contrast ($\Delta\delta = -0.009$, paired-role
bootstrap 95\% CI [-0.238, 0.207]).

\subsection{External / construct / conclusion threats}\label{f.2-external-construct-conclusion-threats}

{%
\begin{longtable}[]{@{}
  >{\raggedright\arraybackslash}p{(\linewidth - 6\tabcolsep) * \real{0.2500}}
  >{\raggedright\arraybackslash}p{(\linewidth - 6\tabcolsep) * \real{0.2500}}
  >{\raggedright\arraybackslash}p{(\linewidth - 6\tabcolsep) * \real{0.2500}}
  >{\raggedright\arraybackslash}p{(\linewidth - 6\tabcolsep) * \real{0.2500}}@{}}
\caption{External / construct / conclusion threats.}\label{tab:p2-28}\\
\toprule\noalign{}
\begin{minipage}[b]{\linewidth}\raggedright
Threat area
\end{minipage} & \begin{minipage}[b]{\linewidth}\raggedright
Threat
\end{minipage} & \begin{minipage}[b]{\linewidth}\raggedright
Class
\end{minipage} & \begin{minipage}[b]{\linewidth}\raggedright
Mitigation
\end{minipage} \\
\midrule\noalign{}
\endfirsthead
\endhead
\bottomrule\noalign{}
\endlastfoot
Representativeness & 12-PUT representativeness & External & follows the shared site selection;
open class principle (cls extensible); scaling to MD/QC/CFD \\
Statistical power & Cross-class statistical power & Conclusion & H3 framed exploratory;
mixed-effects unavailable (Singular at N=60); shifted to class means +
sign test + Friedman + forest plot \\
LLM homogeneity & LLM homogeneity bias & Construct & §C.2 three-LLM rotation +
PUT-class subdivision + 20\% manual sampling; third-party LLM family
validation deferred \\
LLM-source shift & LLM-source distributional shift & External & Hit distribution
DeepSeek 11/15, Claude 4/15, GPT 0/15; the main manuscript's AST-overlap audit argument is categorical
AST-locality, not hit frequency; HP/SI/TF zero-overlap across all 3
sources unaffected \\
Higher-order mutation & HOM equivalence & External & Tool-unreachability claim confined to
first-order syntactic tools; HOM empirical testing listed as future
work \\
\end{longtable}
}

Conclusion threats: multiple comparisons handled per the main-text
policy (per-class Bonferroni $\times$ 4; no cell-level family claimed;
§D.1); N=20 stability handled by 1,000-iteration
bootstrap CI (§D.3). Construct threat, SMS measures epistemological
semantic detection, not production engineering value (out of scope here).

% [thematic break removed]

\section{Degeneration to Classical Mutation Score: Full Proof}\label{g.-sms-ms-degeneration-theorem-full-proof}

\subsection{Notation cross-reference}\label{g.1-notation-cross-reference-with-2.1}

PUT $S_i$, mutation operator family mut (syntactic or semantic),
pooled mutant set $\mathrm{mut}(S_i)$, MR-label set, equivalence tolerance $\varepsilon_{\mathrm{eq}}$, equivalence
sample size $K_{\mathrm{eq}}$, AVP tolerance $\varepsilon_{\mathrm{AVP}}$. Pool
decomposition (three-state refinement, §A.1):
$\mathrm{mut}(S) = \mathrm{certified} \cup \mathrm{killed} \cup \mathrm{survive} \cup \mathrm{unresolved}$
(disjoint). Strict SMS formula (the reported reading; in the
$u = 0$ regime it equals the classical form):

\[
\mathrm{SMS}_{i,k} = \frac{|\mathrm{killed}_{i,k}|}{|\mathrm{killed}_{i,k}| + |\mathrm{survive}_{i,k}|}
= \frac{|\mathrm{killed}_{i,k}|}{|\mathrm{mut}(S_i)| - |\mathrm{equiv}_{i,k}|} \quad (u_{i,k} = 0)
\]

\subsection{Degenerate-limit definition}\label{g.2-degenerate-limit-definition-r-8-p1-3-revision-3-joint-conditions}

The degenerate limit $L = L_{\mathrm{equiv}} \wedge L_{\mathrm{killed}} \wedge L_{\mathrm{mut}}$
consists of three \textbf{joint conditions} (paired axes), each
controlling one layer of the SMS formula. L1-L6 are not 6 independent
axes; the pairing reflects dependency structure across the SMS formula
layers. Equivalently, the same conjunction regroups orthogonally as
$L = L_{\mathrm{lim}} \wedge L_{\mathrm{switch}}$, where
$L_{\mathrm{lim}} = L1 \wedge L2 \wedge L3$ collects the quantitative
limit axes and $L_{\mathrm{switch}} = L4 \wedge L5 \wedge L6$ collects
the discrete configuration switches (reference-anchored identity
relation; rule-based syntactic operators; imperative deterministic PUT
class).

{%
\begin{longtable}[]{@{}
  >{\raggedright\arraybackslash}p{(\linewidth - 4\tabcolsep) * \real{0.3333}}
  >{\raggedright\arraybackslash}p{(\linewidth - 4\tabcolsep) * \real{0.3333}}
  >{\raggedright\arraybackslash}p{(\linewidth - 4\tabcolsep) * \real{0.3333}}@{}}
\caption{Degenerate-limit definition.}\label{tab:p2-29}\\
\toprule\noalign{}
\begin{minipage}[b]{\linewidth}\raggedright
Joint condition
\end{minipage} & \begin{minipage}[b]{\linewidth}\raggedright
Component axes
\end{minipage} & \begin{minipage}[b]{\linewidth}\raggedright
Pairing rationale
\end{minipage} \\
\midrule\noalign{}
\endfirsthead
\endhead
\bottomrule\noalign{}
\endlastfoot
\textbf{$L_{\mathrm{equiv}}$} (Lemma G.1) & L1: $\varepsilon_{\mathrm{eq}} \to 0$; L2: $K_{\mathrm{eq}} \to \infty$ & L1 alone:
equiv stays probabilistic ($K_{\mathrm{eq}}$ cannot cover $\mathcal{D}_P$). L2 alone: bitwise
equality diluted by $\varepsilon_{\mathrm{eq}}$. Both simultaneously degenerate equiv to
classical behavioural equivalence outside the exception set
$\mathcal{N}_{\mathrm{exc}}$. \\
\textbf{$L_{\mathrm{killed}}$} (Lemma G.2) & L3: $\varepsilon_{\mathrm{AVP}}^k \to 0$; L4: MR-label set =
$\{\mathrm{MR}_{\mathrm{eq}}^{S_i}\}$, the reference-anchored identity
relation ($r = \mathrm{id}$, $R(y,y') \equiv y = y'$, follow-up output
compared against the recorded reference output of $S_i$) & L3 alone: $\varepsilon_{\mathrm{AVP}} \to 0$ still allows
non-trivial MR (monotonicity, convergence). L4 alone: equality still
carries $\varepsilon_{\mathrm{AVP}}$ tolerance; without the identity
restriction the comparison $S_i(x)$ vs $s'(r(x))$ retains a metamorphic
follow-up and is not the classical same-input oracle. Both simultaneously degenerate killed to
classical difference detection. \\
\textbf{$L_{\mathrm{mut}}$} (Lemma G.3) & L5: $\mathrm{mut}_j \to$ rule-based syntactic (Mothra
AOR/ROR/SDL/CRP); L6: cls(I) $\subseteq$ \{imperative deterministic\} & L5 alone:
syntactic ops on probabilistic/ML may still trigger domain-semantic
subsets. L6 alone: imperative programs still mutable by OS/HP/TF/SI.
Both simultaneously degenerate mut(S) to Jia \& Harman syntactic mutant
set. \\
\end{longtable}
}

\subsection{Lemmas for the Pool Decomposition}\label{g.3-lemmas-three-state-decomposition-degenerates-under-l}

\textbf{Lemma G.1 (equiv degeneration; exception-set
qualification).} Assume
$\mathrm{supp}(\mathcal{D}_P) \supseteq \mathcal{X}_{\mathrm{adm}}$ and
fix the switch component $L_{\mathrm{switch}}$. Under
$L_{\mathrm{equiv}}$ (L1 $\wedge$ L2), semantic-class equivalence
(E1 $\wedge$ E2) degenerates to classical behavioural equivalence
\textbf{outside an exception set $\mathcal{N}_{\mathrm{exc}}$}: over
floating-point (finite) input domains, $\mathcal{N}_{\mathrm{exc}}$ is
the finite set of pathological inputs (Not-a-Number propagation and
overflow/underflow traps), empty in the exact discrete case; under the
continuum idealisation, $\mathcal{N}_{\mathrm{exc}}$ has
$\mathcal{D}_P$-measure zero.

\emph{Proof.} Under the fixed $L_{\mathrm{switch}}$ the relation set is
the reference-anchored identity relation
$\mathrm{MR}_{\mathrm{eq}}^{S_i}$, whose AVP check of a program $Q$
compares $Q$'s outputs against the recorded reference outputs of $S_i$
on the executed inputs. The E2 limit below yields exact output equality
on its covered domain. Exact equality makes the original and mutant AVP
verdicts identical for every non-negative AVP tolerance, so E1 is
automatically satisfied on the E2 limit. No
$\varepsilon_{\mathrm{AVP}}\to0$ premise is used here; that L3 axis
belongs only to the killed-set degeneration of Lemma G.2. E2
($\|S_i(x) - s'(x)\| < \varepsilon_{\mathrm{eq}}$ for $K_{\mathrm{eq}}$ samples)
degenerates in two regimes. Over floating-point (finite) input domains,
L2 ($K_{\mathrm{eq}} \to \infty$) with full-support sampling
($\mathrm{supp}(\mathcal{D}_P) \supseteq \mathcal{X}_{\mathrm{adm}}$)
covers every admissible input almost surely, and L1
($\varepsilon_{\mathrm{eq}} \to 0$) sharpens the per-input test to exact
equality, giving
$\forall x \in \mathcal{X}_{\mathrm{adm}} \setminus \mathcal{N}_{\mathrm{exc}}: S_i(x) = s'(x)$
with $\mathcal{N}_{\mathrm{exc}}$ the finite pathological set
($\mathcal{N}_{\mathrm{exc}} = \emptyset$ in the exact discrete case).
Under the continuum idealisation, sampled agreement at
$\varepsilon_{\mathrm{eq}} \to 0$ over an almost-surely dense
full-support sample forces equality $\mathcal{D}_P$-almost everywhere,
i.e.\ outside a $\mathcal{D}_P$-measure-zero
$\mathcal{N}_{\mathrm{exc}}$; pointwise equality on all of
$\mathcal{X}_{\mathrm{adm}}$ is not claimed in this regime, and the
degenerate equivalence claim is read on
$\mathcal{X}_{\mathrm{adm}} \setminus \mathcal{N}_{\mathrm{exc}}$
throughout. This matches \citet{d7c38286734419b52de4262c9802ebdfcf4b9447} §3 classical
equivalent-mutant definition up to the exception set.

\emph{Support-assumption counterexample (non-removability).} Let
$\mathcal{X}_{\mathrm{adm}} = [0, 2]$ but
$\mathrm{supp}(\mathcal{D}_P) = [0, 1]$, and let $s'$ differ from $S_i$
only on $(1, 2]$. Every E2 sample agrees at every $K_{\mathrm{eq}}$, so
the limit judgement certifies equivalence, yet $s'$ is behaviourally
non-equivalent on $\mathcal{X}_{\mathrm{adm}}$; the degeneration target
is therefore false without
$\mathrm{supp}(\mathcal{D}_P) \supseteq \mathcal{X}_{\mathrm{adm}}$.

\textbf{Lemma G.2 (killed degeneration).} Under L3 $\wedge$ L4, killed
degenerates to classical same-input difference detection.

\emph{Proof.} L4 restricts the MR-label set to
$\{\mathrm{MR}_{\mathrm{eq}}^{S_i}\}$ with identity source
transformation ($r = \mathrm{id}$), $R(y, y') \equiv y = y'$, and the
follow-up output compared against the recorded reference output of
$S_i$. The killed predicate requires
$\mathrm{AVP}(S_i, mr) = \text{pass}$, which holds identically for the
reference-anchored relation, and $\mathrm{AVP}(s', mr) = \text{fail}$,
which under L3 ($\varepsilon_{\mathrm{AVP}} \to 0$) is exactly
$\exists x: s'(x) \neq S_i(x)$: classical same-input difference
detection against the original program. Without the identity
restriction the comparison $S_i(x)$ vs $s'(r(x))$ retains a metamorphic
follow-up and is not the classical oracle, and without the reference
anchor an identity relation evaluated within a single deterministic
program passes trivially on both $S_i$ and $s'$ and detects nothing;
both clauses of L4 are therefore load-bearing.

\textbf{Lemma G.3 (mut degeneration).} Under L5 $\wedge$ L6, the pooled
$\mathrm{mut}(S_i)$
degenerates to the syntactic mutant set of \citet{d7c38286734419b52de4262c9802ebdfcf4b9447}.

\emph{Proof.} L5 switches every generator family to rule-based syntactic operators (AOR,
ROR, SDL, CRP, UOI, standard Mothra/Proteum sets); L6 restricts PUTs
to imperative deterministic programs, excluding triggering conditions
for semantic operators on probabilistic/ML programs. Under this
configuration, $\mathrm{mut}(S_i)$ is the literature-defined syntactic mutant
set, independent of domain semantics.

\subsection{Degeneration Theorem: Detailed Proof}\label{g.4-theorem-9.1-detailed-proof}

\textbf{Theorem G (SMS $\to$ MS degeneration; stated in the main text
as the degeneration theorem).} Fix the switch component
$L_{\mathrm{switch}} = L4 \wedge L5 \wedge L6$ and assume
$\mathrm{supp}(\mathcal{D}_P) \supseteq \mathcal{X}_{\mathrm{adm}}$.
Along the limit component $L_{\mathrm{lim}} = L1 \wedge L2 \wedge L3$,
outside the exception set $\mathcal{N}_{\mathrm{exc}}$ of Lemma G.1,

\[ \mathrm{SMS}_{i,k} \xrightarrow{L_{\mathrm{lim}}} \mathrm{MS}_{i} := \frac{|\mathrm{killed}_{i}^{\mathrm{classic}}|}{|\mathrm{mut}^{\mathrm{syntax}}(S_i)| - |\mathrm{equiv}_{i}^{\mathrm{classic}}|} \]

\emph{Proof.} Combine Lemmas G.1-G.3:

\begin{itemize}

\item
  Numerator: under L3 $\wedge$ L4, $\mathrm{killed}_{i,k} \to
  \mathrm{killed}_{i}^{\mathrm{classic}}$ (Lemma G.2); L4 makes $\mathrm{MR}_{i,k}$ trivial
  over $k$ (only the reference-anchored identity relation remains), so
  the subscript $k$ degenerates.
\item
  Denominator $|\mathrm{mut}(S_i)|$: under L5 $\wedge$ L6 $\to$
  $|\mathrm{mut}^{\mathrm{syntax}}(S_i)|$ (Lemma G.3).
\item
  Denominator $|\mathrm{equiv}_{i,k}|$: under L1 $\wedge$ L2 with the
  support assumption $\to$
  $|\mathrm{equiv}_{i}^{\mathrm{classic}}|$ (Lemma G.1;
  outside $\mathcal{N}_{\mathrm{exc}}$), and the unresolved set empties
  into the classical dichotomy, so the strict and conservative readings
  coincide in the limit.
\end{itemize}

Substituting into the SMS formula gives $\mathrm{MS}_{i}$; the
restriction to any single generator family is immediate.

\textbf{Corollary G (LRCA trivialisation).} Under $L$, C
= \{C1, \ldots, C5\} degenerates to \{C1\}.

\emph{Sketch.} Each of C2-C5's triggering precondition depends on at
least one $L_j$ being violated. When the three joint conditions
$L_{\mathrm{equiv}} \wedge L_{\mathrm{killed}} \wedge L_{\mathrm{mut}}$
hold simultaneously,
every non-trivial-space dimension (MR non-triviality, AVP tolerance
non-zero, non-empty equiv set, MR-design DOF, class-mapping openness)
closes; C2-C5 triggering set becomes empty (read off from §A.2 decision
tree). The per-C\_k to per-L\_j minimum-sufficient mapping depends on
the main manuscript's LRCA-diagnosis engineering thresholds; we do not claim one-to-one
correspondence at formal level, readers can trace through §A.2 + §C.3.
Under L, $\texttt{suspect\_share} \to 0$, LRCA reports only C1, SMS degenerates to
single-layer metric consistent with \citet{d7c38286734419b52de4262c9802ebdfcf4b9447} MS engineering
structure.

\textbf{Empirical consistency.} Theorem G and the LRCA-trivialisation
corollary jointly
show that the SMS-based empirical conclusions (Cliff's delta in the aligned-versus-cross results,
Friedman chi\^{}2 in the cross-class consistency results, and Spearman rho in the SMS-vs-Pattern-Coverage analysis) are structurally consistent
with existing \citet{d7c38286734419b52de4262c9802ebdfcf4b9447} literature in classical
syntactic-mutation scenarios, no metric-level semantic fragmentation.

\subsection{Independent formal audit of the theory extension}\label{g.5-independent-formal-audit}

The theory-extension results (the detection-window, interval-soundness,
and gap-attribution theorems with their lemmas, corollary, and remarks,
and the repaired degeneration theorem above) were checked by an
independent formal audit against a fixed eight-item checklist: (1)
every theorem premise is defined in the body; (2) no circularity; (3)
the scope of the AVP-determinism assumption of the kill-witness lemma;
(4) consistency of the gap theorem's S5/exact-checker premises with the
$\xi$ reporting rule; (5) satisfiability of the window theorem's
constants and Lipschitz hypothesis per PUT class; (6) the degeneration
theorem's support assumption and exception-set wording; (7) notation
conformance against the frozen symbol registry; (8) line-by-line
checkability of the proofs. The audit raised blocking and minor
findings (premise hygiene, the $\xi$ empty-kill convention and
one-sidedness, the reference-anchored degenerate oracle, estimand
definedness); all were repaired and verified closed in follow-up
rounds, for a final verdict of PASS. The audit protocol and the signed
report are archived with the replication materials.

\subsection{Interval soundness: full proofs}\label{g.6-interval-proofs}

\emph{Proof of the protocol-level kill-witness lemma.} A kill supplies
an \(mr\) for which
\(\mathrm{AVP}(S_i,mr)=\mathrm{pass}\) and
\(\mathrm{AVP}(s',mr)=\mathrm{fail}\). Hence E1 fails on that relation.
The three-state protocol defines E1 or E2 failure as
\texttt{CONFIRMED\_NON\_EQUIVALENT}; no determinism or tolerance-margin
premise is needed. Therefore no killed mutant is unresolved. \(\square\)

\textbf{Scope of the determinism hypothesis (quantitative output
witness).} The
evaluation procedure of a relation fixes its executed input tuple
(seeded input generation is part of the procedure), runs the program on
that tuple, applies the observation $\mathrm{obs}$, and computes the
verdict by a fixed threshold functional at tolerance
$\varepsilon_{\mathrm{AVP}}$. For deterministic PUTs the verdict is a
fixed function $V_{mr}$ of the tuple of observed outputs on the common
executed inputs. For stochastic PUTs the pipeline replaces each AVP call
by a majority vote over $N = 20$ repetitions; the verdict is then a
fixed function of the $N$-repeat aggregated observation, and every
occurrence of ``observed output'' is read at that aggregated level: the
witness produced below is a divergence of aggregated observations, in
the AVP verdict semantics, not a claim about any single random run.
(The LRCA layer's 0.80 fail-ratio cutoff is a separate robustness
annotation on killed mutants and plays no role in the kill verdict.)
The stability hypothesis (clause R2, §G.8) is not removable: with
relation residuals $\tau - \delta_1$ (pass on $S_i$) and
$\tau + \delta_2$ (fail on $s'$), $\delta_1 + \delta_2$ arbitrarily
small, a kill can arise from pointwise output divergence below
$\varepsilon_{\mathrm{eq}}$, and bare determinism cannot yield the
required witness.

\emph{Proof of the quantitative output-witness proposition.} Fix $mr$ with
$\mathrm{AVP}(S_i, mr) = \mathrm{pass}$ and
$\mathrm{AVP}(s', mr) = \mathrm{fail}$. Suppose every executed input $x$
satisfied
$\|\mathrm{obs}(\Phi_{S_i}(x)) - \mathrm{obs}(\Phi_{s'}(x))\| \leq \varepsilon_{\mathrm{eq}}$.
Then the observed-output tuple of $s'$ is a pointwise
$\leq \varepsilon_{\mathrm{eq}}$ perturbation of that of $S_i$, and R2
forces $V_{mr}$ to return the same verdict on both, contradicting
pass/fail. Hence some executed $x$ has divergence exceeding
$\varepsilon_{\mathrm{eq}}$: the divergence witness of the three-state
classification. $\square$

\emph{Proof of the interval theorem.} By the kill-witness lemma the
killed count $\kappa$ lies inside the confirmed non-equivalent
population $n$, so the ground-truth denominator is
$n + u_{\mathrm{neq}}$ with the same numerator $\kappa$.
\emph{(Interval and attainability.)} For fixed $\kappa \geq 0$,
$n \geq 1$ the map $t \mapsto \kappa/(n+t)$ is non-increasing on
$t \geq 0$; applying it to $0 \leq u_{\mathrm{neq}} \leq u$ gives the
sandwich, with both endpoints attained at $u_{\mathrm{neq}} = u$ and
$u_{\mathrm{neq}} = 0$, each consistent with the available evidence.
\emph{(Width.)}
$\kappa/n - \kappa/(n+u) = \kappa u/(n(n+u)) = \mathrm{SMS}^{\mathrm{strict}} \cdot u/(n+u)$.
\emph{(Monotonicity, four motions.)} (1) An equivalence certificate on
an unresolved survivor maps $(n, \kappa, u) \to (n, \kappa, u-1)$: the
upper endpoint is unchanged and the lower endpoint rises, so the
interval weakly narrows (strictly when $\kappa > 0$, $u \geq 1$). (2) A
divergence witness on an unresolved survivor maps
$(n, \kappa, u) \to (n+1, \kappa, u-1)$: the lower endpoint
$\kappa/(n+u)$ is unchanged and the upper endpoint falls to
$\kappa/(n+1)$, so the interval weakly narrows; unresolved mutants are
survivors by the kill-witness lemma, so $\kappa$ is unchanged. (3) An
unresolved survivor killed by a newly added MR acquires the witness,
$(n, \kappa, u) \to (n+1, \kappa+1, u-1)$, an instance of motion (4)
with $\Delta = 0$, $j = 1$. (4) For $R \subseteq R'$ with the
three-state classification held fixed, OR-aggregation over a superset
can only add kills: with $\Delta \geq 0$ new kills among confirmed
survivors and $j \geq 0$ among unresolved survivors, the counts update
to $n' = n + j$, $\kappa' = \kappa + \Delta + j$, $u' = u - j$, and
$\mathrm{SMS}^{\mathrm{strict}\prime} = (\kappa+\Delta+j)/(n+j) \geq (\kappa+j)/(n+j) \geq \kappa/n$
(since $\kappa \leq n$), while
$\mathrm{SMS}^{\mathrm{cons}\prime} = (\kappa+\Delta+j)/(n+u) \geq \kappa/(n+u)$:
both endpoints are non-decreasing. The proviso is not removable: if the
classification were recomputed against $R'$, the E1 quantification
domain would grow and a new relation with differing verdicts in the
non-kill direction would convert an unresolved survivor into a confirmed
one, a motion-(2) event lowering the strict endpoint. Across all four
motions the conservative endpoint never decreases: it is a monotone
lower bound of the ground-truth score under every admissible evidence
update. $\square$

\subsection{Block structure: full proofs}\label{g.7-gap-proofs}

\emph{Proof of the closure lemma (exact checkers are strong MRs).} The
violation-set identity is the exact-checker definition read as a set
equation. Each stratum invariant is a predicate over the observable
behaviour, and the tolerance-indexed satisfaction evaluates that
predicate on the observed behaviour; membership of the violation set
therefore depends on a program only through its observed behaviour, so
observationally equivalent programs violate together: the set is a
union of equivalence classes, i.e.\ closed, and the strong-MR
definition applies. Properness: the union of two distinct strata's
violation classes is closed but pins to neither stratum. $\square$

\textbf{Premise reading.} (i) reads: every admitted non-equivalent
mutant is stratified, violating its unique target stratum (S3 witness)
and satisfying all others (S5), so the fibers $\{M_j\}$ partition
$M_{\mathrm{neq}}$ and $\sum_j w_j = 1$; mutants outside this reading
(active-off-taxonomy, multi-stratum violators) are excluded by the
premise, and their off-diagonal kills are exactly what $\xi$ registers
(the diagnostic is one-sided: premise violations that kill only on the
labelled stratum leave $\xi$ unchanged, so observed $\xi$ lower-bounds
the deviation mass). (ii) attaches to every $r \in R$ a stratum label,
unique by the non-redundancy convention. (iii) imports, for every
$r \in R$, the residual budget and margin-dominance clause R1 of §G.8:
the S5-satisfied stratum contributes zero violation, the budget bounds
the executed residual by $\Delta_r + p\bar\eta$, and R1 places that
below $\varepsilon_{\mathrm{tol}}$.

\emph{Proof of the gap theorem.} \emph{(No cross-stratum flag.)} Let
$r \in R$ be an exact checker of $\psi_l$ and $m \in M_j$ with
$l \neq j$. By (i), $m$ satisfies $\psi_l$; by (iii) the executed
residual of $r$ on $m$ stays below $\varepsilon_{\mathrm{tol}}$, so the
tolerance verdict is pass and $r$ does not flag $m$. The kill predicate
requires a fail on the mutant, so $r$ does not kill it.
\emph{(Uncovered fibers survive.)} If $j \notin \mathrm{Cov}(R)$, every
$r \in R$ has label $l \neq j$, so no $r$ flags any member of $M_j$ and
OR-aggregation yields no kill. \emph{(Block diagonality.)} Indexing rows
by fibers and columns by checker strata, the argument shows every
off-diagonal entry is zero. \emph{(Decomposition.)} The fibers partition
$M_{\mathrm{neq}}$, so
$\mathrm{SMS}(R) = \sum_j w_j\,\mathrm{SMS}_j(R)$ with
$\mathrm{SMS}_j(R) = 0$ for uncovered $j$; hence
$1 - \mathrm{SMS}(R) = \sum_{j \notin \mathrm{Cov}(R)} w_j + \sum_{j \in \mathrm{Cov}(R)} w_j (1 - \mathrm{SMS}_j(R))$,
an exact identity with non-negative terms. \emph{(Computability.)}
$w_j$ from fiber-label counts, $\mathrm{SMS}_j(R)$ from row-wise kill
rates, $\mathrm{Cov}(R)$ from the checker labels carried by the kill
matrix's column index: no equivalence oracle beyond the
$M_{\mathrm{neq}}$ certification and no ground truth beyond the labels.
$\square$

\emph{Proof of the cross-zero corollary.} If no covered stratum has
positive mass,
$\mathrm{SMS}(R) = \sum_{j \in \mathrm{Cov}(R)} w_j\,\mathrm{SMS}_j(R) = 0$.
$\square$

\emph{Justification of the identifiability remark.} Same-stratum
membership of a killed mutant's signature is the block-diagonal matrix
reread row-wise: a killed member of $M_j$ can only be flagged by
checkers with label $j$, so its nonempty signature certifies the kill
and names the fiber; with one exact checker per covered stratum the
killed subpopulations separate every covered stratum. Survivor rows are
identically zero, hence observationally indistinguishable: the
below-window remainder of every covered fiber (window theorem, clause
(ii)) and every uncovered fiber produce the same all-pass observation,
so survivor attribution is carried by generation-time labels, not by
execution outcomes. $\square$

\subsection{Detection window: hypotheses, regime, and full proofs}\label{g.8-window-proofs}

\textbf{Hypotheses.} \emph{(H-a, additive residual budget.)} On the
executed tuple at input $x \in D_r$, the residual the checker computes
for the mutant satisfies
$|\varepsilon_r(x; m) - \varepsilon_{\mathrm{viol}}(m)| \leq \Delta_r + p\bar\eta$:
the ideal mutant residual decomposes into the structural violation
contribution plus the structure-preservation residual, bounded by
$\Delta_r$ because the edit perturbs the declared structure, not the
discretisation; and the executed tuple aggregates $p$ program
executions, each observation carrying noise at most $\bar\eta$. The
pairwise relations that dominate the catalogue have $p = 2$; the
convergence-order relations evaluate $p = 4$ grid resolutions and their
regression verdict additionally carries a design-dependent conditioning
constant, so their window bounds are read with the $p = 4$ budget.
\emph{(H-b, Lipschitz transfer.)} The violation functional is
$L_r$-Lipschitz in $\varepsilon_{\mathrm{viol}}(m)$; H-b is used only
for the dose-response transfer between the nominal-parameter axis and
the realized violation axis, not for bounds (i)--(iii), which are
stated on the realized axis. Estimability of $L_r$ is assessed per PUT
class: linear and quasi-linear kernels admit closed-form or regression
estimates, while chaotic and stochastic-structural kernels do not admit
a stable $L_r$ on the raw observable, and any dose-response consumption
of H-b is restricted to the estimable subset and pre-registered as
such. \emph{(H-c, non-degenerate-margin regime.)} Two clauses: (R1)
margin dominance,
$\mu_r = \varepsilon_{\mathrm{tol}} - \Delta_r > p\bar\eta$ strictly;
and (R2) $\varepsilon_{\mathrm{eq}}$-separation, no executed residual
within $\varepsilon_{\mathrm{eq}}$ of the threshold, making the verdict
invariant under pointwise obs-output perturbations of magnitude at most
$\varepsilon_{\mathrm{eq}}$. R1 is consumed by (i) and (iii) and by the
gap theorem's premise (iii); R2 is consumed by the kill-witness lemma.
For unbounded noise the bound $|\eta| \leq \bar\eta$ is a
concentration-style model assumption and stochastic-case statements
hold at the confidence attached to $c$. \emph{(H-d, magnitude
realization.)} The executed tuple realizes the declared magnitude; for
uniform parametric violation templates every admissible input realizes
it, and unrealized magnitudes only weaken (i), reverting detection to
finite-input search, while (ii) and (iii) are unaffected.

\emph{Proof of (i).} By H-a's lower side,
$\varepsilon_r(x; m) \geq \varepsilon_{\mathrm{viol}}(m) - \Delta_r - p\bar\eta > \varepsilon_{\mathrm{tol}}$,
so the checker flags the mutant (strict exceedance); by R1 every
executed residual of the original program is at most
$\Delta_r + p\bar\eta < \varepsilon_{\mathrm{tol}}$, so the checker
passes the original; the kill predicate's two conjuncts hold. $\square$

\emph{Proof of (ii).} By H-a's upper side, on every executed tuple
$\varepsilon_r(x; m) \leq \varepsilon_{\mathrm{viol}}(m) + \Delta_r + p\bar\eta < \varepsilon_{\mathrm{tol}}$,
so the checker never flags the mutant and the kill predicate's second
conjunct fails; the verdict on the original is not needed. $\square$

\emph{Proof of (iii).} The lower edge is the contrapositive of (ii); at
equality the executed residual is at most $\varepsilon_{\mathrm{tol}}$,
which under strict exceedance still yields no flag, so killed magnitudes
exceed
$\varepsilon_{\mathrm{tol}} - \Delta_r - p\bar\eta$
strictly. Magnitudes at or beyond $\varepsilon_{\mathrm{crash}}$ violate
S4 latency and are excluded from the admitted universe by the crash
oracle, so the kill region lies within
$(\varepsilon_{\mathrm{tol}} - \Delta_r - p\bar\eta,\ \varepsilon_{\mathrm{crash}})$.
$\square$

\emph{Proof of the weak-MR remark.} $\mu_r < 0$ iff
$\Delta_r > \varepsilon_{\mathrm{tol}}$; since $\Delta_r$ is a supremum
over $D_r$, inputs exist whose correct-program residual exceeds
$\varepsilon_{\mathrm{tol}}$ (no attainment needed). If validation
executes an input whose residual exceeds
$\varepsilon_{\mathrm{tol}} + p\bar\eta$, the observed residual stays
above $\varepsilon_{\mathrm{tol}}$ under any admissible noise and the
flag is guaranteed; inputs at that level exist exactly when
$\mu_r < -p\bar\eta$. In the band
$(\varepsilon_{\mathrm{tol}}, \varepsilon_{\mathrm{tol}} + p\bar\eta]$
the flag occurs unless noise masks it, and in all cases validation
sampling must reach such inputs at all. A flagged correct program
breaks the kill predicate's pass-on-original conjunct, so the relation
contributes no kills and is rejected upstream at MR validation.
$\square$

\emph{Proof of the stochastic remark.} For a pairwise relation
($p = 2$), substituting
$\bar\eta = c\,\sigma_{\mathrm{out}}/\sqrt{N} + \eta_{\mathrm{det}}$
into (i) at target magnitude $\varepsilon^\dagger$ and solving for $N$
gives
$\sqrt{N} > 2c\,\sigma_{\mathrm{out}}/(\varepsilon^\dagger - \varepsilon_{\mathrm{tol}} - \Delta_r - 2\eta_{\mathrm{det}})$,
provided the denominator is positive; squaring gives the stated bound,
which is sufficient with strict inequality; necessity holds under the
boundary-attaining noise model used to derive the budget. If the side
condition fails, no repeat budget guarantees detection and the remedy
is a different observable. $\square$

\textbf{Structure-fate correspondence.} $\Delta_r$ places each MR on a
structure-preservation axis: $\Delta_r = 0$, exact preservation;
$0 < \Delta_r \leq \varepsilon_{\mathrm{tol}}$, approximate
preservation (the strong/tolerance MR of the main text);
$\Delta_r > \varepsilon_{\mathrm{tol}}$, structure broken (weak MR, the
first remark's $\mu_r < 0$); $\Delta_r(h) \to 0$ under refinement,
asymptotic preservation. The window theorem prices the same $\Delta_r$
into two-sided detection bounds: how much violation magnitude a valid
relation must convert into a kill and how much it must miss.

\section{Adjoint Extension Study}\label{h.-adjoint-extension-arm}

This appendix gives the full specification of the confirmatory adjoint
extension arm summarised in the main manuscript. The arm
exercises the adjoint invariant $\psi_6$ on three third-party
linear-operator libraries, disjoint from the 12 PUTs and the
pre-registered 60-cell. Each library carries a real fixed defect in its
adjoint code path as a ground-truth anchor:

\begin{itemize}

\item
  \textbf{scipy.sparse.linalg LinearOperator} (linear algebra; BSD-3).
  Program under test: the scaled-operator adjoint $M^{\mathrm{H}}$ for
  $M = \alpha A$ with complex scale $\alpha$. Real anchor: scipy issue
  \#8900 / PR \#8962 (the adjoint dropped the conjugate on $\alpha$,
  returning $\alpha A^{\mathrm{H}}$ instead of $\bar{\alpha}
  A^{\mathrm{H}}$).
\item
  \textbf{pylops.signalprocessing.FFT} with \texttt{real=True} (signal
  processing; BSD-3). Program under test: the real-FFT adjoint
  \texttt{\_rmatvec}; the adjoint identity is the library's own
  \emph{dot-test}. Real anchor: pylops issue \#268 / PR \#292 (commit
  \texttt{bd3a919}), a wrong scaling of the middle/Nyquist frequency bin
  that made the adjoint inconsistent with the forward.
\item
  \textbf{jax} \texttt{linear\_transpose} / \texttt{vjp} (automatic
  differentiation; Apache-2.0). Program under test: the autodiff-derived
  transpose $M^{\mathrm{T}}$ of a linear operator containing an inverse
  real FFT, $M(x) = \mathrm{irfft}(\mathrm{rfft}(x)\cdot w)$. Real anchor:
  jax issue \#6223 / PR \#6232 (commit \texttt{112ae41}), an
  \texttt{irfft} transpose rule that reversed order with an off-by-one
  index, so \texttt{vjp} silently produced a non-adjoint transpose.
\end{itemize}

For each subject the MR families are adj.self (self-adjointness) and
adj.dual (scaled/dual-operator transpose). The mutant pool spans
adjoint-breaking fault mechanisms (drop-conjugate, sign flip, scale and
phase errors, transpose-without-conjugation, Nyquist mis-scaling), with
the real fixed defect included among them, plus semantically equivalent
mutants (scalar-multiplication reorder, conjugation of a real adjoint
output) that a strong MR must \emph{not} kill. Because the pre-fix
releases are uninstallable on the host platform (no \texttt{arm64} wheel),
each real defect is reproduced in extracted-diff form on the current
release, mirroring the boundary study; $\varepsilon_{\mathrm{tol}}$ is an
explicit factor throughout.

{%
\begin{longtable}[]{@{}p{2.0cm}p{3.2cm}p{2.5cm}p{1.8cm}p{1.8cm}@{}}
\caption{Adjoint extension study, forward kill.}\label{tab:suppl-adjoint}\\
\toprule\noalign{}
Substrate & Operator class (real anchor) & Active killed / total & Equiv.\ survived & Forward SMS \\
\midrule\noalign{}
\endfirsthead
\endhead
\bottomrule\noalign{}
\endlastfoot
scipy LinearOperator & linear algebra (\#8900) & 5 / 5 & 2 / 2 & \textbf{1.00} \\
pylops FFT real & signal processing (\#292) & 4 / 4 & 2 / 2 & \textbf{1.00} \\
jax vjp / transpose & automatic differentiation (\#6223) & 4 / 4 & 2 / 2 & \textbf{1.00} \\
\end{longtable}
}

Across all three substrates the correct program survives, every
non-equivalent (active) mutant is killed, including the three real
fixed defects, and no semantically equivalent mutant is killed, so the
forward SMS is 1.00 with zero false positives on the equivalent set. The
scipy \#8900 defect is shared with the strong-boundary study in the main manuscript (adj.self case)
and is not counted as independent corroboration twice; the pylops \#292
and jax \#6223 defects are unique to this arm. The pools are
hand-constructed and low-cardinality, so the arm demonstrates the duality
forward rather than measuring a high-adequacy MR set on a discovered fault
population (main-paper adjoint-arm reproduction-fidelity threat).

\section{Real-Defect Evidence Summary}\label{i.-result-level-real-defect-evidence-ledger}

This appendix records only the result-level facts used by the main
manuscript's real-defect evidence slice. It deliberately excludes
candidate mining, rejected cases, repository tooling, and benchmark-release
narrative. The purpose is auditability for the semantic-mutation argument:
kill-rate, semantic-stratum alignment, and real-defect detection are related
but distinct constructs.

{%
\small
\begin{longtable}[]{@{}p{0.18\linewidth}p{0.52\linewidth}p{0.24\linewidth}@{}}
\caption{Result-level real-defect evidence summary.}\label{tab:realdefect-ledger}\\
\toprule\noalign{}
Audit item & Evidence content used in the manuscript & Main-text role \\
\midrule\noalign{}
\endfirsthead
\endhead
\bottomrule\noalign{}
\endlastfoot
Verified stratum counts & 34 verified MR-detectable defects across five
meta-patterns: \(m_{\mathrm{mono}}\) 10, \(m_{\mathrm{conv}}\) 9,
\(m_{\mathrm{inv}}\) 6, \(m_{\mathrm{adj}}\) 5, and
\(m_{\mathrm{rev}}\) 4 & Shows that the declared semantic strata are not
only toy labels \\
Mutation-phase aggregate & The semantic-mutation track exceeds the
closest baseline by +0.101 on mean score (Holm-adjusted \(p = 0.046\),
Cliff's \(\delta = +0.247\)); the remaining pairwise contrasts are not
significant & Supports a modest, qualified
mutation-score advantage rather than a dominance claim \\
Real-defect face & The semantic-mutation track detects all 34 tabled
real defects; the two generic baseline groups have nonzero detection in
7/34 and 8/34 cases, and the two ablated variants miss 19/34 and 17/34
cases, respectively, with 11 shared misses & Separates real-defect detection from aggregate mutant
kill-rate \\
Non-nesting cases & Four verified non-nesting defects (A-LAPACK-004,
A-OPENBLAS-001, B-POCKETFFT-002, and E-ORDINARYDIFFEQ-001) spanning
linear
algebra, FFT, and ordinary-differential-equation solvers & Supports the semantic-effect fiber view rather
than a single scalar hierarchy \\
\end{longtable}
\normalsize
}

This summary is intentionally a result table, not a corpus or benchmark
description. The main manuscript uses it only to support construct
separation: a relation family can score similarly on aggregate mutant
kill-rate while differing sharply on real-defect detection, and different
relation families can observe non-nested semantic-effect fibers.

\bibliographystyle{ACM-Reference-Format}
\bibliography{references}

\end{document}
